\SourcePage{117}
\SGAchapter{XII}{对射影代数概形的应用}{Applications aux schémas algébriques projectifs}
\SGAsection{1}{射影对偶定理与有限性定理}
\footnote{本讲写于1963年1月，比1962年6月的口头报告详细得多。}
下面的定理实质上已包含在 FAC\footnote{J.-P. Serre，« Faisceaux algébriques cohérents »，\emph{Ann. of Math.} \textbf{61} (1955)，197--278页。} 中（不过当时 Serre 还没有\Term{module}{模}\Term{sheaf}{层}的 Ext 语言\footnote{编者注：喜欢数学史的读者会饶有兴趣地读到 Grothendieck 于1955年12月15日写给 Serre 的信，以及后者同年12月22日的回信；见 Pierre Colmez 和 Jean-Pierre Serre 编，\emph{Correspondance Grothendieck--Serre}，Documents Mathématiques，第2卷，Société Mathématique de France，Paris，2001。}），它是\Term{local-duality}{局部对偶}定理（第 IV 讲）的整体类似物；后者正是以它为蓝本的。
\begin{SGAresult}{定理}{1.1}
设 $k$ 是\Term{field}{域}，$X=\mathbf P_k^r$ 是 $k$ 上的 $r$ 维\Term{projective-space}{射影空间}，$F$ 是 $X$ 上变动的\Term{coherent}{凝聚}\Term{module}{模}。则有关于 $F$ 的 $\partial$-\Term{functor}{函子}的\Term{isomorphism}{同构}：
\[
 H^i(X,F)'\xleftarrow{\sim}\operatorname{Ext}^{r-i}(X;F,\Omega^r_{X/k}),\tag{1}
\]
其中约定
\[
 \Omega^r_{X/k}=\mathcal O_{\mathbf P_k^r}(-r-1).\tag{2}
\]
\end{SGAresult}
\textbf{注记。}当然，这个\Term{module}{模}也是 $X$ 在 $k$ 上的 $r$ 次\Term{relative-differentials}{相对微分}\Term{module}{模}。在这种形式下，若 $X$ 是 $k$ 上\Term{proper}{固有}且\Term{smooth}{光滑}的\Term{scheme}{概形}，定理仍然成立（\Term{projective}{射影}情形见 A. Grothendieck，« Théorèmes de dualité pour les faisceaux algébriques cohérents »，Séminaire Bourbaki，1957年5月）\footnote{更一般的\Term{duality}{对偶}定理，参见第 IV 讲末尾引用的 Hartshorne 讨论班。}。当 $F$ \Term{locally-free}{局部自由}时，就得到 Serre \Term{duality}{对偶}定理
\[
 H^i(X,F)'\xrightarrow{\sim}H^{r-i}\bigl(\mathcal Hom_{\mathcal O_X}(F,\Omega^r_{X/k})\bigr).
\]
定理1.1对我们的目的已经足够（它也重新给出了上述文献中 $X$ 在 $k$ 上\Term{projective}{射影}的情形）。

\SourcePage{118}
\Term{homomorphism}{同态}(1)由 Yoneda \Term{pairing}{配对}
\[
 H^i(X,F)\times\operatorname{Ext}^{r-i}(X;F,\Omega^r_{X/k})\longrightarrow H^r(X,\Omega^r_{X/k})\tag{3}
\]
和一个熟知的\Term{isomorphism}{同构}（参见 FAC 或 EGA III 2.1.12）导出：
\[
 H^r(X,\Omega^r_{X/k})=H^r(\mathbf P_k^r,\mathcal O_{\mathbf P_k^r}(-r-1))\xrightarrow{\sim}k.\tag{4}
\]
证明(1)是\Term{isomorphism}{同构}的方法同\Term{local-duality}{局部对偶}定理的情形一样：注意 $H^r(X,F)$ 作为 $F$ 的\Term{functor}{函子}是\Term{right-exact}{右正合}的（因为 $n>r$ 时 $H^n(X,F)=0$），而且每个\Term{coherent}{凝聚}\Term{module}{模}都\Term{isomorphism}{同构}于一些形如 $\mathcal O(-m)$、$m$ 充分大的\Term{module}{模}的\Term{direct-sum}{直和}的一个商。于是对 $i$ 作\Term{descent}{下降}归纳，问题便归结为对形如 $\mathcal O(-m)$ 的\Term{sheaf}{层}进行验证，而这包含在熟知的显式计算中（FAC 或 EGA III 2.1.12）。并且可以假设 $-m\leq-r-1$，此时 $i\neq r$ 就有 $H^i(X,\mathcal O(-m))=0$。
\begin{SGAresult}{推论}{1.2}
给定\Term{coherent}{凝聚}\Term{module}{模} $F$，当 $m$ 充分大时有\Term{canonical}{典范}\Term{isomorphism}{同构}
\[
 H^i(X,F(-m))'\xrightarrow{\sim}H^0\bigl(X,\mathcal Ext^{r-i}_{\mathcal O_X}(F,\Omega^r_{X/k})(m)\bigr),\tag{5}
\]
其中 $'$ 表示\Term{duality}{对偶}\Term{vector-space}{向量空间}。
\end{SGAresult}
事实上，在\Term{projective-space}{射影空间} $X$ 上，对任意两个\Term{coherent}{凝聚}\Term{sheaf}{层} $F,G$ 以及充分大的 $n$，有\Term{canonical}{典范}\Term{isomorphism}{同构}：
\[
 \operatorname{Ext}^n(X;F(-m),G)\simeq\operatorname{Ext}^n(X;F,G(m))
 \xrightarrow{\sim}H^0\bigl(X,\mathcal Ext^n_{\mathcal O_X}(F,G)(m)\bigr).\tag{6}
\]
前两项的\Term{isomorphism}{同构}显然对任意 $m$ 都成立；上述结论来自整体 Ext 的\Term{spectral-sequence}{谱序列}
\[
 H^p\bigl(X,\mathcal Ext^q_{\mathcal O_X}(F,G(m))\bigr)\Longrightarrow\operatorname{Ext}^{\bullet}(X;F,G(m)).
\]
当 $m$ 充分大时，由于
\[
 \mathcal Ext^q_{\mathcal O_X}(F,G(m))\simeq\mathcal Ext^q_{\mathcal O_X}(F,G)(m)
\]
且 $\mathcal Ext^q_{\mathcal O_X}(F,G)$ 是\Term{coherent}{凝聚}\Term{sheaf}{层}，这个\Term{spectral-sequence}{谱序列}\Term{degeneration}{退化}。所以(5)由(6)与(1)得到。
\begin{SGAresult}{推论}{1.3}
给定 $i,F$，下列条件\Term{equivalence}{等价}：
\begin{enumerate}
\item[(i)] 当 $m$ 充分大时，$H^i(X,F(-m))=0$。
\item[(i bis)] $H^i(X,F(\cdot))=\bigoplus_{m\in\mathbf Z}H^i(X,F(-m))$ 是\Term{finite-type}{有限型} $S$-\Term{module}{模}，其中 $S=k[t_0,\ldots,t_r]$。
\item[(ii)] $\mathcal Ext^{r-i}_{\mathcal O_X}(F,\Omega^r_{X/k})=0$。
\item[(ii bis)] $\mathcal Ext^{r-i}_{\mathcal O_X}(F,\mathcal O_X)=0$。
\item[(iii)] 对 $X$ 的每个\Term{closed-point}{闭点} $x$，$H^i_x(F_x)=0$。
\item[(iv)] 对 $X$ 的\Term{punctured-cone}{去顶点投射锥} $\widetilde X=\operatorname{Spec}(S)-\operatorname{Spec}(k)$ 的每个\Term{closed-point}{闭点} $x$，$H^{i+1}_x(\widetilde F_x)=0$；这里 $\widetilde F$ 表示 $F$ 沿\Term{canonical}{典范}\Term{morphism}{态射} $\widetilde X\to X$ 的\Term{inverse-image}{逆像}。
\end{enumerate}
\end{SGAresult}
\SourcePage{119}
\begin{proof}
(i)$\Leftrightarrow$(i bis)，因为 $H^i(X,F(\cdot))$ 中由次数 $\geq\nu$ 的\Term{homogeneous-component}{齐次分量}的\Term{direct-sum}{直和}组成的子\Term{module}{模}是\Term{finite-type}{有限型} $S$-\Term{module}{模}（事实上当 $i\neq0$ 时，它甚至是\Term{finite-type}{有限型} $k$-\Term{module}{模}），参见 FAC 或 EGA III 2.2.1、2.3.2。

(i)$\Leftrightarrow$(ii) 来自推论1.2。

(ii)$\Leftrightarrow$(ii bis)，因为 $\Omega^r_{X/k}$ \Term{local-scheme}{局部}\Term{isomorphism}{同构}于 $\mathcal O_X$。

(ii bis)$\Leftrightarrow$(iii) 来自 $\mathcal O_{X,x}$ 的\Term{local-duality}{局部对偶}定理（它是\Term{dimension}{维数}为 $r$ 的\Term{regular-local-ring}{正则局部环}）；根据这个定理，$\mathcal Ext^{r-i}_{\mathcal O_X}(F,\mathcal O_X)_x$ 的“\Term{duality}{对偶}”可与 $H^i_x(F_x)$ 等同（V 2.1）。

(ii bis)\Term{equivalence}{等价}于类似的关系
\[
 \mathcal Ext^{r-i}_{\mathcal O_{\widetilde X}}(\widetilde F,\mathcal O_{\widetilde X})=0
\]
（因为 $\widetilde X\to X$ \Term{faithfully-flat}{忠实平坦}，所以 $\mathcal Ext^{r-i}_{\mathcal O_X}(F,\mathcal O_X)$ 的\Term{inverse-image}{逆像}\Term{isomorphism}{同构}于 $\mathcal Ext^{r-i}_{\mathcal O_{\widetilde X}}(\widetilde F,\mathcal O_{\widetilde X})$），而最后这个关系由\Term{local-ring}{局部环} $\mathcal O_{X,x}$ 的\Term{local-duality}{局部对偶}定理\Term{equivalence}{等价}于(iv)；这个\Term{local-ring}{局部环}\Term{regular}{正则}，\Term{dimension}{维数}为 $r+1$。
\end{proof}
特别地，把这个结果应用于所有 $i\leq n$，便得到：
\begin{SGAresult}{推论}{1.4}
给定 $n,F$，下列条件\Term{equivalence}{等价}：
\begin{enumerate}
\item[(i)] 当 $i\leq n$ 且 $m$ 充分大时，$H^i(X,F(-m))=0$。
\item[(i bis)] 当 $i\leq n$ 时，$H^i(X,F(\cdot))$ 是\Term{finite-type}{有限型} $S$-\Term{module}{模}。
\item[(ii)] 对 $X$ 的每个\Term{closed-point}{闭点} $x$，$\operatorname{prof}(F_x)>n$。
\item[(iii)] 对 $\widetilde X$ 的每个\Term{closed-point}{闭点} $x$，$\operatorname{prof}(\widetilde F_x)>n+1$。
\end{enumerate}
\end{SGAresult}
推论1.3和1.4的意义在于把整体条件(i)或(i bis)表达成\Term{local-scheme}{局部}条件，即 $H^i_x(X,F_x)$ 或 $H^i_x(X,\widetilde F_x)$ 这样的\Term{local-scheme}{局部}\Term{invariants}{不变量}的消失，或关于\Term{depth}{深度}的不等式。在这种形式下，这些结果对任意\Term{projective}{射影}\Term{scheme}{概形} $X$ 及其上的\Term{very-ample}{甚丰沛}\Term{invertible}{可逆}\Term{sheaf}{层} $\mathcal O_X(1)$ 仍然显然成立：只要用适当的\Term{projective}{射影}\Term{immersion}{浸入} $X\to\mathbf P_k^r$ 来诱导后者即可看出。（当然，在这种一般情形下，1.3(ii)和1.3(ii bis)不再与其他条件\Term{equivalence}{等价}，除非例如假设 $X$ \Term{regular}{正则}。）还可以如下推广到\Term{projective}{射影}\Term{morphism}{态射} $X\to S$ 的情形。
\begin{SGAresult}{命题}{1.5}
设 $f:X\to S$ 是\Term{projective}{射影}\Term{morphism}{态射}，$S$ \Term{noetherian}{诺特}，$\mathcal O_X(1)$ 是 $X$ 上相对于 $S$ \Term{very-ample}{甚丰沛}的\Term{invertible}{可逆}\Term{module}{模}，$F$ 是 $X$ 上相对于 $S$ \Term{flat}{平坦}的\Term{coherent}{凝聚}\Term{module}{模}，$s\in S$，$X_s$ 是 $X$ 在 $s$ 处的\Term{fibre}{纤维}（看成 $k(s)$ 上的\Term{projective}{射影}\Term{scheme}{概形}），$F_s$ 是 $F$ 在 $X_s$ 上诱导的\Term{sheaf}{层}，最后 $i$ 是整数。假设对 $X_s$ 的每个\Term{closed-point}{闭点} $x$，都有 $H^i_x(F_{s,x})=0$（例如 $\operatorname{prof}(F_{s,x})>i$）。则 $s$ 有一个\Term{open-neighborhood}{开邻域} $U$，使同一条件对 $s'\in U$ 成立。此外，对这样的 $U$，有
\[
 R^if_*(F(-m))=0\qquad(m\text{ 充分大}),
\]
\SourcePage{120}
并且，若 $\mathscr S$ 是 $S$ 上由 $\mathscr S^1$ 生成的\Term{quasi-coherent}{拟凝聚}\Term{graded-algebra}{分次代数}，且由
\[
 X=\operatorname{Proj}(\mathscr S),\qquad \mathcal O_X(1)=\operatorname{Proj}(\mathscr S(1))
\]
定义 $X$ 及 $\mathcal O_X(1)$，则 $S$-\Term{module}{模}
\[
 R^if_*(F(\cdot))=\bigoplus_{m\in\mathbf Z}R^if_*(F(m))
\]
在 $U$ 上\Term{finite-type}{有限型}。
\end{SGAresult}
\begin{proof}
把 $X$ 嵌入 $X'=\mathbf P_S^r$，使 $\mathcal O_X(1)$ 由 $\mathcal O_{X'}(1)$ 诱导（把 $S$ 换成 $s$ 的一个\Term{affine}{仿射}邻域后即可做到）。对每个整数 $j$ 及每个 $t\in S$，令\footnote{把纯\Term{local-scheme}{局部}的命题 EGA IV 12.3.4 应用于上述 $E^j$，立刻可得1.5的第一部分，从而避开下面证明的大部分。}
\[
 E^j(t)=\mathcal Ext^j_{\mathcal O_{X'_t}}(F'_t,\mathcal O_{X'_t}(-r-1)).\tag{7}
\]
于是 $E^j(t)$ 是 $X_t$ 上的\Term{coherent}{凝聚}\Term{module}{模}。我断言，当 $t$ 变动时，这族\Term{module}{模}在以下意义下是“\Term{constructible}{可构造}的”：对每个 $t\in S$，\Term{closure}{闭包} $\overline{\{t\}}$（赋予诱导的\Term{reduced}{约化}结构）中存在\Term{nonempty}{非空}\Term{open-subset}{开集} $V$，以及 $X_V=X\times_SV$ 上相对于 $V$ \Term{flat}{平坦}的\Term{coherent}{凝聚}\Term{module}{模} $E^j(V)$，使对任意 $t'\in V$，$E^j(t)$ \Term{isomorphism}{同构}于 $E^j(V)$ 在 $X_t$ 上诱导的\Term{module}{模}。为验证此断言，令 $Z=\overline{\{t\}}$ 并赋予诱导结构，考虑\Term{coherent}{凝聚}\Term{module}{模}
\[
 E^j(Z)=\mathcal Ext^j_{\mathcal O_{X'_Z}}(F_Z,\mathcal O_{X'_Z}(-r-1))
\]
（下标 $Z$ 仍表示在 $Z$ 上作诱导），再取 $Z$ 的\Term{nonempty}{非空}\Term{open-subset}{开集} $V$，使\Term{module}{模} $E^j(Z)$ 在 $V$ 上\Term{flat}{平坦}。这是可能的，因为立刻可验证 $j\notin[0,r]$ 时 $E^j(Z)=0$，因而可以应用 SGA 1 IV 6.11。取 $E^j(V)=E^j(Z)|_{X_V}$，易验证它满足要求。

由上述注记，通过\Term{noetherian}{诺特}归纳，可得 $S$ 的一个\Term{finite}{有限}划分，由形如上述 $V=V(t)$ 的\Term{set}{集合} $V_\alpha$ 组成。对 $E^j(V_i)$ 应用 Serre 定理 EGA III 2.2.1，便知存在整数 $m_0$，使
\[
 R^if_{V_\alpha*}(E^j(V_\alpha))=0\qquad(i\neq0,\ m\geq m_0,\ \text{任意 }j).
\]
由 $E^j(V_\alpha)$ 相对于 $V_\alpha$ 的\Term{flatness}{平坦性}及容易得到的 Künneth 型关系（参见 EGA III 第7节），可得
\[
 H^i(X_t,E^j(t)(m))=0\qquad(i\neq0,\ m\geq m_0,\ \text{任意 }j),
\]
对每个 $t\in V_\alpha$ 成立，因而对每个 $t$ 成立，因为 $V_\alpha$ \Term{covering}{覆盖} $S$。由此和整体 Ext 的\Term{spectral-sequence}{谱序列}，借助1.1，并同1.2的证明一样，得到\Term{isomorphism}{同构}
\[
 H^i(X_t,F_t(-m))'\xleftarrow{\sim}H^0(X_t,E^{r-i}(t)(m))\qquad(m\geq m_0),\tag{8}
\]
其中 $i$ 为任意整数，$t\in S$ 任意。
\SourcePage{121}
现在使用关于 $F_s$ 的假设，它写成
\[
 E^{r-i}(s)=0,\tag{9}
\]
由(8)，\Term{equivalence}{等价}于
\[
 H^i(X_s,F_s(-m))=0\qquad(m\geq m_0).\tag{10}
\]
由于 $F$，从而 $F(-m)$ 相对于 $S$ \Term{flat}{平坦}，由已援用的 Künneth 型关系可知：（对给定 $m$）用 $s$ 附近的 $t$ 替代 $s$，同一关系(10)仍成立，特别对 $s$ 的任意\Term{generalization}{一般化点} $t$ 成立。根据(8)，对这样的一个\Term{generalization}{一般化点}便有
\[
 E^{r-i}(t)=0.\tag{11}
\]
而满足这个关系的 $t\in s$ 所成的\Term{set}{集合}显然\Term{constructible}{可构造}（因为它在每个 $V_\alpha$ 上诱导一个\Term{open-subset}{开集}）；它既然包含 $s$ 的所有\Term{generalization}{一般化点}，便包含 $s$ 的一个\Term{open-neighborhood}{开邻域} $U$。这就证明了1.5的第一条断言。此外，对 $t\in U$，由(11)和(8)推出
\[
 H^i(X_t,F_t(-m))=0\qquad(m\geq m_0,\ t\in U).\tag{12}
\]
根据 Künneth 型关系，这蕴含（事实上比它强得多）
\[
 R^if_*(F(-m))=0\quad\text{在 }U\text{ 上成立，}\quad m\geq m_0.\tag{13}
\]
这证明了1.5的第二条断言。最后一条则立即得到，只要如1.3证明的开头那样进行即可。
\end{proof}
\textbf{注记 1.6。}\footnote{第112页的脚注使本注记更为明确。}如果已经假设对 $s\in S$ 所作的假设在所有 $s'\in S$ 都成立，证明就显著简化（排除了关于\Term{constructibility}{可构造性}的一切考虑）。事实上，当假设 $F_s$ 在 $X$ 的\Term{closed-point}{闭点}处\Term{depth}{深度} $>i$ 时，有一个在 $X$ 上\Term{local-scheme}{局部}的一般命题：把 $X$ 换成\Term{fibre}{纤维} $X_s$ 的适当\Term{open-neighborhood}{开邻域}后，所有 $X_t$ 都满足同一条件（换言之，$X$ 中一个由 $F$ 在\Term{fibre}{纤维}上诱导的\Term{module}{模}所满足的条件定义的\Term{subset}{子集}是\Term{open-subset}{开集}，参见 EGA IV）。这里 $f$ \Term{proper}{固有}，因此可以取形如 $f^{-1}(U)$ 的邻域，从而不经繁琐的逐步\Term{resolution}{分解}便得到1.5的第一条断言。在这个一般情形下，仍可以用上述文献的方法证明：1.5的第一条断言（这里利用 $X$ 在 $S$ 上\Term{projective}{射影}而以整体方法证明）仍来自一个在 $X$ 上纯\Term{local-scheme}{局部}的命题（读者如认为有用，可以明确写出）。

\SourcePage{122}
\SGAsection{2}{射影态射的 Lefschetz 理论：Grauert 比较定理}
这就是下面的定理：
\begin{SGAresult}{定理}{2.1}
设 $f:X\to S$ 是\Term{projective}{射影}\Term{morphism}{态射}，$S$ \Term{noetherian}{诺特}，$\mathcal O_X(1)$ 是 $X$ 上相对于 $S$ \Term{ample}{丰沛}的\Term{invertible}{可逆}\Term{module}{模}，$Y$ 是 $\mathcal O_X(1)$ 的\Term{section}{截面} $t$ 的零点\Term{prescheme}{预概形}，$\mathcal J$ 是定义 $Y$ 的\Term{ideal}{理想}，$X_n$ 是 $\mathcal J^{n+1}$ 定义的 $X$ 的子\Term{prescheme}{预概形}，$\widehat X$ 是 $X$ 沿 $Y$ 的\Term{formal-completion}{形式完备化}，$\widehat f:\widehat X\to S$ 是复合\Term{morphism}{态射} $\widehat X\to X\to S$，$F$ 是 $X$ 上相对于 $S$ \Term{flat}{平坦}的\Term{coherent}{凝聚}\Term{module}{模}。再假设对每个 $s\in S$，$F$ 在\Term{fibre}{纤维} $X_s$ 上诱导的\Term{module}{模} $F_s$ 在该\Term{fibre}{纤维}的各点处\Term{depth}{深度} $>n$，并且 $t$ 是 $F$-\Term{regular}{正则}的。在这些条件下：
\begin{enumerate}
\item[(i)] \Term{canonical}{典范}\Term{homomorphism}{同态}
\[
 R^if_*(F)\longrightarrow R^i\widehat f_*(\widehat F)
\]
当 $i<n$ 时是\Term{isomorphism}{同构}，当 $i=n$ 时是\Term{monomorphism}{单态}。
\item[(ii)] \Term{canonical}{典范}\Term{homomorphism}{同态}
\[
 R^i\widehat f_*(\widehat F)\longrightarrow\varprojlim_m R^if_*(F_m)
\]
当 $i\leq n$ 时是\Term{isomorphism}{同构}。
\end{enumerate}
\end{SGAresult}
\begin{proof}
立刻归结到 $S$ \Term{affine}{仿射}的情形，此时只须证明下述推论。
\end{proof}
\begin{SGAresult}{推论}{2.2}
在2.1的条件下，再假设 $S$ \Term{affine}{仿射}。则：
\begin{enumerate}
\item[(i)] \Term{canonical}{典范}\Term{homomorphism}{同态} $H^i(X,F)\to H^i(\widehat X,\widehat F)$ 当 $i<n$ 时是\Term{isomorphism}{同构}，当 $i=n$ 时是\Term{monomorphism}{单态}。
\item[(ii)] \Term{canonical}{典范}\Term{homomorphism}{同态}
\[
 H^i(\widehat X,\widehat F)\longrightarrow\varprojlim_m H^i(X_m,F_m)
\]
当 $i\leq n$ 时是\Term{isomorphism}{同构}。
\end{enumerate}
\end{SGAresult}
\begin{proof}
把 $\mathcal O_X(1)$ 换成一个\Term{tensor-power}{张量幂}，把 $t$ 换成它的一个幂，便可假设 $\mathcal O_X(1)$ 相对于 $S$ \Term{very-ample}{甚丰沛}。另一方面，$t$ 从而 $t^m$ 是 $F$-\Term{regular}{正则}的，因此乘以 $t^m$ 作为 $F(-m)$ 到 $F$ 的\Term{homomorphism}{同态}是\Term{injective-map}{单射}，所以对每个 $m\geq0$，有\Term{exact-sequence}{正合列}：
\[
 0\longrightarrow F(-m)\xrightarrow{t^m}F\longrightarrow F_m\longrightarrow0.\tag{14}
\]
从而有\Term{cohomology}{上同调}\Term{exact-sequence}{正合列}
\[
 H^i(X,F(-m))\longrightarrow H^i(X,F)\longrightarrow H^i(X,F_m)\longrightarrow H^{i+1}(X,F(-m)).
\]
由1.5，当 $i\leq n$ 且 $m$ 充分大时 $H^i(X,F(-m))=0$，这证明了下述引理。
\end{proof}
\SourcePage{123}
\begin{SGAresult}{引理}{2.3}
当 $m$ 充分大时，\Term{canonical}{典范}\Term{homomorphism}{同态}
\[
 H^i(X,F)\longrightarrow H^i(X,F_m)
\]
在 $i<n$ 时是\Term{bijective-map}{双射}，在 $i=n$ 时是\Term{injective-map}{单射}。
\end{SGAresult}
这表明当 $i<n$ 时，\Term{projective-system}{投射系统} $(H^i(X_m,F_m))_{m\geq0}$ \Term{essentially-constant}{本质上为常系统}，因而当然满足 \Term{mittag-leffler}{Mittag-Leffler 条件}；所以由 EGA $0_{\mathrm{III}}$ 13.3，并考虑到 $\widehat F=\varprojlim F_m$，得到(ii)。另一方面，再考虑到2.3，(i)便显然随之成立。
\begin{SGAresult}{推论}{2.4}
\footnote{编者注：参见 Raynaud 女士的文章中的推论 I.1.4：Raynaud M.，« Théorèmes de Lefschetz en cohomologie des faisceaux cohérents et en cohomologie étale. Application au groupe fondamental »，\emph{Ann. Sci. Éc. Norm. Sup.} (4) \textbf{7} (1974)，29--52页。}
设 $f:X\to S$ 是\Term{projective}{射影}且\Term{flat}{平坦}的\Term{morphism}{态射}，$S$ \Term{local-scheme}{局部}\Term{noetherian}{诺特}，$\mathcal O_X(1)$ 是 $X$ 上相对于 $S$ \Term{ample}{丰沛}的\Term{invertible}{可逆}\Term{module}{模}，$t$ 是此\Term{module}{模}的一个 $\mathcal O_X$-\Term{regular}{正则}\Term{section}{截面}，$Y$ 是 $t$ 的零点子\Term{prescheme}{预概形}，$\widehat X$ 是 $X$ 沿 $Y$ 的\Term{formal-completion}{形式完备化}。假设对每个 $s\in S$，$X_s$ 在其\Term{closed-point}{闭点}处\Term{depth}{深度} $\geq1$（分别地，\Term{depth}{深度} $\geq2$）。则对 $Y$ 的每个\Term{open-neighborhood}{开邻域} $U$，\Term{functor}{函子}
\[
 F\longmapsto\widehat F
\]
从 $U$ 上\Term{locally-free}{局部自由}\Term{coherent}{凝聚}\Term{module}{模}的\Term{category}{范畴}到 $\widehat X$ 上\Term{locally-free}{局部自由}\Term{coherent}{凝聚}\Term{module}{模}的\Term{category}{范畴}是\Term{faithful}{忠实}的（分别地，\Term{fully-faithful}{完全忠实}，即满足 X 2 的 \Term{lefschetz-condition}{Lefschetz 条件} $(\mathrm{Lef})$）。
\end{SGAresult}
\begin{proof}
对 $U$ 上两个\Term{locally-free}{局部自由}\Term{module}{模} $F,G$，引入\Term{module}{模}
\[
 H=\mathcal Hom_{\mathcal O_U}(F,G),
\]
问题就归结为证明\Term{canonical}{典范}\Term{homomorphism}{同态}
\[
 H^0(U,H)\longrightarrow H^0(\widehat U,\widehat H)\tag{15}
\]
是\Term{injective-map}{单射}（分别地，\Term{bijective-map}{双射}）。\Term{module}{模} $H_t$ 在 $X_t$ 的\Term{closed-point}{闭点}处\Term{depth}{深度} $\geq1$（分别地，$\geq2$），因此可以应用2.1，它在 $U=X$ 的情形蕴含2.4的结论。对任意 $U$，注意问题在 $S$ 上是\Term{local-scheme}{局部}的，所以可假设 $S$ \Term{affine}{仿射}。于是 $X$ 上的每个\Term{coherent}{凝聚}\Term{module}{模}都是一个\Term{locally-free}{局部自由}\Term{coherent}{凝聚}\Term{module}{模}的商（因为 $\mathcal O_X(1)$ 是 $X$ 上相对\Term{ample}{丰沛}的\Term{invertible}{可逆}\Term{module}{模}）。由于\Term{duality}{对偶}\Term{module}{模} $H'=\mathcal Hom(H,\mathcal O_U)$ 可以延拓成 $X$ 上的\Term{coherent}{凝聚}\Term{module}{模}，因而\Term{isomorphism}{同构}于 $X$ 上\Term{locally-free}{局部自由}\Term{module}{模}间的某个\Term{homomorphism}{同态}的\Term{cokernel}{余核}，经转置便知可以找到 $X$ 上\Term{locally-free}{局部自由}\Term{module}{模}间的\Term{homomorphism}{同态}
\[
 u':L'^0\longrightarrow L'^1,
\]
\SourcePage{124}
它在 $U$ 上诱导\Term{locally-free}{局部自由}\Term{module}{模}间的\Term{homomorphism}{同态}
\[
 u:L^0\longrightarrow L^1,
\]
使下列序列\Term{exact}{正合}：
\[
 0\longrightarrow H\longrightarrow L^0\xrightarrow{u}L^1.
\]
用五引理（这里化为三引理）及 $H^0$ \Term{functor}{函子}的\Term{left-exact}{左正合}性，归结为证明：在(15)中以 $L^0,L^1$ 代替 $H$ 时，该\Term{homomorphism}{同态}是\Term{injective-map}{单射}（分别地，\Term{bijective-map}{双射}）；这又归结到 $H$ 由 $X$ 上的\Term{locally-free}{局部自由}\Term{module}{模} $H'$ 诱导的情形。而且在非“分别地”的情形，这一步归约甚至不必要，因为(15)的核无论如何都是 $H$ 在 $U$ 上那些在 $Y$ 的某个适当\Term{open-neighborhood}{开邻域} $V$ 上消失的\Term{section}{截面}；而限制\Term{homomorphism}{同态} $H^0(U,H)\to H^0(V,H)$ 是\Term{injective-map}{单射}，因为 $H$ 在 $X$ 中每个不与 $Y$ 相交的\Term{closed-subset}{闭子集} $Z$ 的各点处\Term{depth}{深度} $\geq1$（见下面的引理）。在“分别地”的情形，问题归结为证明
\[
 H^0(X,H')\longrightarrow H^0(U,H')
\]
为\Term{bijective-map}{双射}；这是因为 $H'$ 在不与 $Y$ 相交的闭集 $Z=X-U$ 的每个点处\Term{depth}{深度} $\geq2$。因此只须证明下述引理。
\end{proof}
\begin{SGAresult}{引理}{2.5}
设 $F$ 是 $X$ 上相对于 $S$ \Term{flat}{平坦}的\Term{coherent}{凝聚}\Term{module}{模}，且对每个 $s\in S$，$F_s$ 在 $X_s$ 的每个\Term{closed-point}{闭点}处\Term{depth}{深度} $\geq n$。则对 $X$ 中每个不与 $Y$ 相交的\Term{closed-subset}{闭子集} $Z$，$F$ 在 $Z$ 的所有点处\Term{depth}{深度} $\geq n$。
\end{SGAresult}
\begin{proof}
事实上，对每个 $x\in X$，令 $s=f(x)$，有
\[
 \operatorname{prof}(F_x)\geq\operatorname{prof}(F_{s,x}).\tag{16}
\]
这是因为，把 $\operatorname r(\mathcal O_{X_s,x})$ 中的一条极大 $F_{s,x}$-\Term{regular-sequence}{正则序列}任意\Term{lifting}{提升}，由 SGA 1 IV 5.7 便得到一条 $F_x$-\Term{regular-sequence}{正则序列}。若 $x$ 属于引理2.5中的 $Z$，则 $x$ 必为 $X_s$ 中的\Term{closed-point}{闭点}；换言之，$Z$ 在 $S$ 上\Term{finite}{有限}。事实上，$Z$（赋予 $X$ 诱导的结构）作为\Term{projective}{射影}的 $X$ 的闭子\Term{prescheme}{预概形}，在 $S$ 上\Term{projective}{射影}；又作为 $X-Y$ 的闭子\Term{prescheme}{预概形}，在 $S$ 上\Term{affine}{仿射}。
\end{proof}
\textbf{注记 2.6。}假设对每个 $s\in S$，$t$ 在 $\mathcal O_{X_s}(1)$ 上诱导的\Term{section}{截面} $t_s$ 是 $\mathcal O_{X_s}$-\Term{regular}{正则}的（由 SGA 1 IV 5.7，这蕴含 $t$ 是 $\mathcal O_X$-\Term{regular}{正则}的）。则所作的假设在\Term{base-extension}{基底扩张} $S'\to S$（$S'$ \Term{local-scheme}{局部}\Term{noetherian}{诺特}）下\Term{stable}{稳定}。所以结论在任意\Term{base-change}{基变换}后仍然\Term{effective}{有效}。

\SourcePage{125}
\SGAsection{3}{射影态射的 Lefschetz 理论：存在定理}
\begin{SGAresult}{定理}{3.1}
\footnote{编者注：不假设\Term{flatness}{平坦性}的版本，见 Raynaud M.，« Théorèmes de Lefschetz en cohomologie des faisceaux cohérents et en cohomologie étale. Application au groupe fondamental »，\emph{Ann. Sci. Éc. Norm. Sup.} (4) \textbf{7} (1974)，29--52页，定理 II.3.3。}
设 $f:X\to S$ 是\Term{projective}{射影}\Term{morphism}{态射}，$S$ \Term{noetherian}{诺特}，$\mathcal O_X(1)$ 是 $X$ 上相对于 $S$ \Term{ample}{丰沛}的\Term{invertible}{可逆}\Term{module}{模}，$X_0$ 是 $\mathcal O_X(1)$ 的\Term{section}{截面} $t$ 的零点子\Term{prescheme}{预概形}，$\widehat X$ 是 $X$ 沿 $X_0$ 的\Term{formal-completion}{形式完备化}，$\mathfrak F$ 是 $\widehat X$ 上的\Term{coherent}{凝聚}\Term{module}{模}，$F_0$ 是它在 $X_0$ 上诱导的\Term{module}{模}。再假设：
\begin{enumerate}
\item[a)] $\mathfrak F$ 相对于 $S$ \Term{flat}{平坦}。
\item[b)] 对每个 $s\in S$，$t$ 在 $X_s$ 上诱导的\Term{section}{截面} $t_s$ 是 $\mathfrak F_s$-\Term{regular}{正则}的（这蕴含 $F_0$ 也相对于 $S$ \Term{flat}{平坦}，参见 SGA 1 IV 5.7）。
\item[c)] 对每个 $s\in S$，$F_{0s}$ 在 $X_{0s}$ 的\Term{closed-point}{闭点}处\Term{depth}{深度} $\geq2$。
\end{enumerate}
再假设 $S$ 有一个\Term{ample}{丰沛}\Term{invertible}{可逆}\Term{sheaf}{层}。在这些条件下，存在 $X$ 上的\Term{coherent}{凝聚}\Term{module}{模} $F$，以及其\Term{formal-completion}{形式完备化} $\widehat F$ 与 $\mathfrak F$ 的一个\Term{isomorphism}{同构}。
\end{SGAresult}
这个命题将由下述结果导出。
\begin{SGAresult}{推论}{3.2}
在上述条件 a)、b)、c) 下，有：
\begin{enumerate}
\item[(i)] $S$ 上的\Term{module}{模} $\widehat f_*(\mathfrak F)$ 是\Term{coherent}{凝聚}的，因而对每个 $n$，$S$ 上的\Term{module}{模} $\widehat f_*(\mathfrak F(n))$ 是\Term{coherent}{凝聚}的。
\item[(ii)] 当 $n$ 充分大时，\Term{canonical}{典范}\Term{homomorphism}{同态} $\widehat f^*\widehat f_*(\mathfrak F(n))\to\mathfrak F(n)$ 是\Term{surjective-map}{满射}。
\end{enumerate}
\end{SGAresult}
\begin{proof}[定理3.1的证明]
暂且承认这个推论，来证明3.1。由3.1中的最后一个假设，把 $\mathcal O_X(1),t$ 换成适当的幂后，可归结到 $X=\mathbf P_S^r$ 的情形。我断言，还可以假设对每个 $s$ 都有 $t_s\neq0$。否则由 b) 就有 $\mathfrak F_s=0$；或由 Nakayama 引理\Term{equivalence}{等价}地，$F_{0s}=0$，即 $s$ 不属于 $\operatorname{Supp}F_0$ 在 $f$ 诱导的\Term{morphism}{态射} $f_0:X_0\to S$ 下的像。由 a)、b)，$F_0$ 相对于 $S$ \Term{flat}{平坦}，故这个像 $S'$ 是开的；显然只须在限制到 $S'$ 上所得到的情形中证明3.1的结论，因为这样得到的 $X|_{S'}$ 上的\Term{coherent}{凝聚}\Term{module}{模} $F'$ 是 $X$ 上某个\Term{coherent}{凝聚}\Term{module}{模} $F$ 的限制，而这个 $F$ 就满足要求。所以除假设 a)、b)、c) 外，还可以假设下列条件成立：
\begin{enumerate}
\item[a$'$)] $\mathcal O_X$ 相对于 $S$ \Term{flat}{平坦}。
\item[b$'$)] 对每个 $s\in S$，\Term{section}{截面} $t_s$ 是 $\mathcal O_{X_s}$-\Term{regular}{正则}的。
\item[c$'$)] 对每个 $s\in S$，$\mathcal O_{X_{0s}}$ 在 $X_{0s}$ 的\Term{closed-point}{闭点}处\Term{depth}{深度} $\geq2$。
\end{enumerate}
（只须取 $X=\mathbf P_S^r$，其中 $r\geq3$，这是允许的。）
\SourcePage{126}
由3.2，可以找到\Term{epimorphism}{满态}
\[
 \widehat L\longrightarrow\mathfrak F\longrightarrow0,\tag{17}
\]
其中 $L$ 是 $X$ 上形如 $f^*(G)(-n)$ 的\Term{module}{模}，$G$ 是 $S$ 上的\Term{locally-free}{局部自由}\Term{coherent}{凝聚}\Term{module}{模}。事实上，只须对充分大的 $n$ 把 $S$ 上的\Term{coherent}{凝聚}\Term{module}{模} $\widehat f_*(\mathfrak F)$ 表成这样的 $G$ 的商。另一方面，关于 $f,t$ 的假设 a)、b)、c) 蕴含 $\widehat L$ 满足与 $\mathfrak F$ 相同的条件 a)、b)、c)。容易推出(17)的核也是如此，因此可以对核应用同样的论证，把 $\mathfrak F$ 表成\Term{homomorphism}{同态}
\[
 \widehat L'\longrightarrow\widehat L\tag{18}
\]
的\Term{cokernel}{余核}，其中 $L,L'$ 是 $X$ 上的\Term{locally-free}{局部自由}\Term{module}{模}。由 a$'$)、c$'$)的第二部分以及2.1或2.4，\Term{homomorphism}{同态}(18)来自 $X$ 上\Term{module}{模}间的\Term{homomorphism}{同态} $L'\to L$。现在只须取 $F$ 为 $L'\to L$ 的\Term{cokernel}{余核}，便得所求。
\end{proof}
还须证明3.2。在讨论班上，这曾用一种颇为繁琐的办法完成：把一切解释为 $X$ 相对于 $S$ 的\Term{punctured-cone}{去顶点投射锥}上的\Term{cohomology}{上同调}，从而归结到定理2.1。现在看来，下面的方法更直接，也更令人满意（虽然实质上相同）。它在于注意：在第 IX 讲第2节（并沿用该讲记号），\Term{morphism}{态射} $f:\mathfrak X\to\mathfrak X'$ 为\Term{adic-morphism}{进态射}的假设，在通过 EGA $0_{\mathrm{III}}$ 13.7.7 证明2.1时根本未被使用；只须改为假设 $\mathfrak X$ 也为\Term{adic-formal-scheme}{进形式概形}，选择 $\mathfrak X'$ 的定义\Term{ideal}{理想} $\mathcal J$ 和 $\mathfrak X$ 的定义\Term{ideal}{理想} $\mathcal I$，使 $\mathcal J\mathcal O_{\mathfrak X}\subset\mathcal I$，定义 $\mathscr S=\operatorname{gr}_{\mathcal J}(\mathcal O_{\mathfrak X'})$，并考虑 $\operatorname{gr}_{\mathcal I}(\mathfrak F)$。无论如何，2.1都可直接应用于本节的\Term{morphism}{态射} $\widehat f:\widehat X\to S$，只须取 $\mathcal J=0$。因此，要验证 $\widehat f_*(\mathfrak F)$ \Term{coherent}{凝聚}，根据上述文献，只须验证 $i=0,1$ 时 $R^if_{0*}(\operatorname{gr}_{\mathcal I}(\mathfrak F))$ 在 $S$ 上\Term{coherent}{凝聚}。为此注意，由 a)、b)，所考虑的\Term{module}{模}就是
\[
 \bigoplus_{m\geq0}R^if_{0*}(F_0(-m)),
\]
而由假设 c) 和1.5，它确实\Term{coherent}{凝聚}。

这证明了3.2(i)。对于3.2(ii)，需要下面的引理。
\begin{SGAresult}{引理}{3.3}
在3.1的条件 a)、b)、c) 下，令
\[
 G_m=\widehat f_*(\mathfrak F(\cdot)_m)=\bigoplus_n\widehat f_*(\mathfrak F_m(n)).
\]
则\Term{projective-system}{投射系统} $(G_m)$ 满足 \Term{mittag-leffler}{Mittag-Leffler 条件}。
\end{SGAresult}
\begin{proof}
可以假设 $S$ \Term{affine}{仿射}，环为 $A$。设 $\mathscr S$ 是正\Term{graded}{分次}的\Term{finite-type}{有限型} $A$-\Term{graded-algebra}{分次代数}，$t'\in\mathscr S_1$，使 $X$ \Term{immersion}{浸入} $\operatorname{Proj}(\mathscr S)$，其中 $\mathcal O_X(1)$ 由 $\mathcal O(1)$ 诱导，\Term{section}{截面} $t$ 是 $t'$ 的像。给 $\mathscr S$ 赋予 $\mathcal J$-进\Term{filtration}{滤过}，其中 $\mathcal J=t'\mathscr S$，
\SourcePage{127}
并在 $X_0$ 上的阿贝尔\Term{sheaf}{层}\Term{category}{范畴}中考虑 $\mathfrak F(\cdot)_m$ 的\Term{projective-system}{投射系统}。仍然满足 EGA $0_{\mathrm{III}}$ 13.7.7\footnote{按第 IX 讲第85页所述修正。} 的预备条件，而且 $i=0,1$ 时，$H^i(X_0,\operatorname{gr}(\mathfrak F(\cdot)))$ 是\Term{finite-type}{有限型} $\operatorname{gr}_{\mathcal J}(\mathscr S)$-\Term{module}{模}。事实上，由于 $t'$ 是 $\mathfrak F$-\Term{regular}{正则}的，立刻看出，作为 $(\mathscr S/t'\mathscr S)[T]$ 上的\Term{module}{模}（$\operatorname{gr}_{\mathcal J}(\mathscr S)$ 是这个环的商），所考虑的\Term{module}{模}可等同于
\[
 H^i(X_0,F_0(\cdot))\otimes_{\mathscr S/t'\mathscr S}(\mathscr S/t'\mathscr S)[T].
\]
根据1.5，当 $i=0,1$ 时，$H^i(X_0,F_0(\cdot))$ 是\Term{finite-type}{有限型} $\mathscr S$-\Term{module}{模}，从而是\Term{finite-type}{有限型} $\mathscr S/t'\mathscr S$-\Term{module}{模}，这证明了我们的断言。因此可以在 $n=1$ 时应用 $0_{\mathrm{III}}$ 13.7.7，便证明了3.3。
\end{proof}
得到这一点后（仍假设 $S$ \Term{affine}{仿射}，这对证明3.2(ii)是允许的），取 $m_0$，使 $m\geq m_0$ 蕴含
\[
 \operatorname{Im}(G_m\to G_0)=\operatorname{Im}(G_{m_0}\to G_0),
\]
于是两边也等于
\[
 \operatorname{Im}(\varprojlim G_m\to G_0)
 =\operatorname{Im}\bigl(\widehat f_*(\mathfrak F(\cdot))\to f_*F_0(\cdot)\bigr).
\]
现在注意，当 $n$ 充分大时，$\mathfrak F_{m_0}(n)$ 由它的\Term{section}{截面}生成，所以 $F_0(n)$ 由可\Term{lifting}{提升}到 $\mathfrak F_{m_0}$ 的\Term{section}{截面}生成，因而（由 $m_0$ 的选择）由可\Term{lifting}{提升}到 $\mathfrak F$ 的\Term{section}{截面}生成。所以 $\mathfrak F$ 的\Term{section}{截面}生成 $F_0$，再由 Nakayama 引理也生成 $\mathfrak F$。这就证明了3.2(ii)，从而证明了3.1。
\begin{SGAresult}{推论}{3.4}
设 $f:X\to S$ 是\Term{projective}{射影}且\Term{flat}{平坦}的\Term{morphism}{态射}，$S$ \Term{local-scheme}{局部}\Term{noetherian}{诺特}，$\mathcal O_X(1)$ 是 $X$ 上相对于 $S$ \Term{ample}{丰沛}的\Term{invertible}{可逆}\Term{module}{模}，$t$ 是这个\Term{module}{模}的\Term{section}{截面}，且对每个 $s\in S$，它在\Term{fibre}{纤维} $X_s$ 上诱导的\Term{section}{截面} $t_s$ 是 $\mathcal O_{X_s}$-\Term{regular}{正则}的；$X_0$ 是 $t$ 的零点子\Term{prescheme}{预概形}，$\widehat X$ 是 $X$ 沿 $X_0$ 的\Term{formal-completion}{形式完备化}。假设对每个 $s\in S$，$X_{0s}$ 在其\Term{closed-point}{闭点}处\Term{depth}{深度} $\geq2$（即 $X_s$ 在 $X_{0s}$ 的\Term{closed-point}{闭点}处\Term{depth}{深度} $\geq3$），且 $X_s$ 在其\Term{closed-point}{闭点}处\Term{depth}{深度} $\geq2$。在这些条件下，偶对 $(X,X_0)$ 满足第 X 讲第2节的\Term{effective-lefschetz-condition}{有效 Lefschetz 条件} $(\mathrm{Leff})$，即：
\begin{enumerate}
\item[a)] 对 $X$ 中 $X_0$ 的每个\Term{open-neighborhood}{开邻域} $U$，\Term{functor}{函子} $F\mapsto\widehat F$ 从 $U$ 上\Term{locally-free}{局部自由}\Term{coherent}{凝聚}\Term{module}{模}的\Term{category}{范畴}到 $\widehat X$ 上\Term{locally-free}{局部自由}\Term{coherent}{凝聚}\Term{module}{模}的\Term{category}{范畴}\Term{fully-faithful}{完全忠实}。
\item[b)] 对 $\widehat X$ 上每个\Term{locally-free}{局部自由}\Term{coherent}{凝聚}\Term{module}{模} $\mathfrak F$，存在 $X_0$ 的一个\Term{open-neighborhood}{开邻域} $U$ 以及 $U$ 上的\Term{locally-free}{局部自由}\Term{coherent}{凝聚}\Term{module}{模} $F$，使 $\mathfrak F$ \Term{isomorphism}{同构}于 $\widehat F$。
\end{enumerate}
\end{SGAresult}
\begin{proof}
a)已在2.4中于更弱的条件下指出。对于 b)，应用3.1便得到结论，至少当 $S$ \Term{noetherian}{诺特}且具有绝对\Term{ample}{丰沛}的\Term{invertible}{可逆}\Term{module}{模}时如此，特别当 $S$ \Term{affine}{仿射}时如此。事实上，若 $F$ 是 $X$ 上的\Term{coherent}{凝聚}\Term{module}{模}，且 $\widehat F$ \Term{isomorphism}{同构}于 $\mathfrak F$，因而\Term{locally-free}{局部自由}，则 $F$ 在 $X_0$ 的某个邻域 $U$ 上\Term{locally-free}{局部自由}，而 $F|_U$ 就满足所需条件。但现在注意，由2.5，这样的 $F$ 在\Term{immersion}{浸入} $U\to X$ 下的像是\Term{coherent}{凝聚}的，而且不依赖于所选的解 $(U,F)$
\SourcePage{128}
（因为由 a)，两个解在 $X_0$ 的邻域内一致）。确切地说，可以找到 $X$ 上的\Term{coherent}{凝聚}\Term{module}{模} $F$ 以及\Term{isomorphism}{同构} $\widehat F\xrightarrow{\sim}\mathfrak F$，使 $F$ 在 $X$ 中那些在其\Term{fibre}{纤维}内闭合的点处\Term{depth}{深度} $\geq2$；这在唯一\Term{isomorphism}{同构}的意义下确定了 $F$。由这个唯一性性质，在 $S$ 的各\Term{affine}{仿射}\Term{open-subset}{开集}上诱导后找到的解可以\Term{gluing}{粘合}，从而得到整个 $X$ 上的\Term{coherent}{凝聚}\Term{module}{模} $F$ 和\Term{isomorphism}{同构} $\widehat F\simeq\mathfrak F$。把 $F$ 限制到它为\Term{free-module}{自由模}的点所成的\Term{open-subset}{开集} $U$，便得到所求。
\end{proof}
根据2.4和3.4，在\Term{projective}{射影}\Term{algebraic-scheme}{代数概形}及其一个“\Term{hyperplane-section}{超平面截面}”的情形中，可以运用第 X、XI 讲中关于条件 $(\mathrm{Lef})$ 和 $(\mathrm{Leff})$ 建立的一般事实。例如：
\begin{SGAresult}{推论}{3.5}
设 $X$ 是\Term{projective}{射影}\Term{algebraic-scheme}{代数概形}，赋有\Term{ample}{丰沛}\Term{invertible}{可逆}\Term{module}{模} $\mathcal O_X(1)$；$t$ 是这个\Term{module}{模}的一个 $\mathcal O_X$-\Term{regular}{正则}\Term{section}{截面}，$X_0$ 是 $t$ 的零点子\Term{scheme}{概形}。假设 $X$ 在其\Term{closed-point}{闭点}处\Term{depth}{深度} $\geq2$（分别地，并且在 $X_0$ 的\Term{closed-point}{闭点}处\Term{depth}{深度} $\geq3$）。则 $\pi_0(X_0)\to\pi_0(X)$ 是\Term{bijective-map}{双射}；特别地，$X$ \Term{connected}{连通}当且仅当 $X_0$ \Term{connected}{连通}。并且，在 $X_0$ 中选定一个\Term{geometric-basepoint}{几何基点}后，$\pi_1(X_0)\to\pi_1(X)$ 是\Term{surjective-map}{满射}；更一般地，对每个\Term{open-subset}{开集} $U\supset X_0$，\Term{homomorphism}{同态} $\pi_1(X_0)\to\pi_1(U)$ 是\Term{surjective-map}{满射}（分别地，\Term{homomorphism}{同态} $\pi_1(X_0)\to\varprojlim_U\pi_1(U)$ 是\Term{bijective-map}{双射}）。在“分别地”的情形，若进一步假设 $X$ 中每个不属于 $X_0$ 的\Term{closed-point}{闭点}的\Term{local-ring}{局部环}是\Term{pure}{纯}的（3.2）（例如\Term{regular}{正则}，或仅为\Term{complete-intersection}{完全交}），则 $\pi_1(X_0)\to\pi_1(X)$ 是\Term{isomorphism}{同构}。
\end{SGAresult}
\begin{proof}
应用2.4和3.4。注意，在“分别地”的情形，$X$ 在 $X_0$ 的\Term{closed-point}{闭点}处\Term{depth}{深度} $\geq3$ 这一假设蕴含 $X$ 的所有\Term{dimension}{维数} $\neq0$ 的\Term{irreducible-component}{不可约分支}的\Term{dimension}{维数}都 $\geq3$（注意这样的分支必与 $X_0$ 相交，再考察交中的一个\Term{closed-point}{闭点}即可看出）。
\end{proof}
\textbf{注记。}当 $X$ \Term{normal}{正规}且在各点处\Term{dimension}{维数} $\geq2$ 时，它在其\Term{closed-point}{闭点}处\Term{depth}{深度} $\geq2$，因而满足3.5中非“分别地”的条件。在这种情形下，可以借助 Bertini 定理更初等地证明 $\pi_1(X_0)\to\pi_1(X)$ 的\Term{surjective-map}{满射}性（参见 SGA 1 X.2.10）。若再假设 $X_0$ \Term{normal}{正规}，且 $X$ 在各点处\Term{dimension}{维数} $\geq3$，则满足3.5中“分别地”的条件。在这种情形下，3.5由 Grauert 建立（事实上，由于\Term{normality}{正规性}假设，借助一些技巧可以绕开\Term{existence-theorem}{存在定理}3.1）。Grauert 的这个证明就是本讨论班所研究的“Lefschetz 理论”的出发点。

\SourcePage{129}
\begin{SGAresult}{推论}{3.6}
设 $X,\mathcal O_X(1),t,X_0$ 如3.5。假设 $X$ 在其\Term{closed-point}{闭点}处\Term{depth}{深度} $\geq2$，并且
\[
 H^i(X_0,\mathcal O_{X_0}(-n))=0
\]
对 $n>0$ 和 $i=1$（分别地，对 $i=1$ 和 $i=2$）成立；由1.4，这蕴含 $X_0$ 在其\Term{closed-point}{闭点}处\Term{depth}{深度} $\geq2$（分别地，$\geq3$），即 $X$ 在 $X_0$ 的\Term{closed-point}{闭点}处\Term{depth}{深度} $\geq3$（分别地，$\geq4$）。在这些条件下，对 $X_0$ 的每个\Term{open-neighborhood}{开邻域} $U$，$\operatorname{Pic}(U)\to\operatorname{Pic}(X_0)$ 是\Term{injective-map}{单射}；特别地，$\operatorname{Pic}(X)\to\operatorname{Pic}(X_0)$ 是\Term{injective-map}{单射}（分别地，$\varprojlim_U\operatorname{Pic}(U)\to\operatorname{Pic}(X_0)$ 是\Term{bijective-map}{双射}）。在“分别地”的情形，若再假设 $X$ 中每个不属于 $X_0$ 的\Term{closed-point}{闭点}的\Term{local-ring}{局部环}是\Term{parafactorial}{副阶乘}的（3.1）（例如\Term{regular}{正则}，或更一般地为\Term{complete-intersection}{完全交}），则 $\operatorname{Pic}(X)\to\operatorname{Pic}(X_0)$ 是\Term{bijective-map}{双射}。
\end{SGAresult}
\begin{proof}
应用 XI 3.12和3.13，并注意“分别地”的假设蕴含 $X$ 的\Term{dimension}{维数} $\neq0$ 的\Term{irreducible-component}{不可约分支}的\Term{dimension}{维数}都 $\geq4$。特别把这个结果应用于\Term{projective-space}{射影空间}中\Term{dimension}{维数} $\geq4$ 的整体\Term{complete-intersection}{完全交} $X$，得到：
\end{proof}
\begin{SGAresult}{推论}{3.7}
设 $X$ 是\Term{dimension}{维数} $\geq3$ 的\Term{algebraic-scheme}{代数概形}，且为某个\Term{scheme}{概形} $\mathbf P_k^r$ 中的\Term{complete-intersection}{完全交}。则 $\operatorname{Pic}(X)$ 是由\Term{sheaf}{层} $\mathcal O_X(1)$ 的类生成的\Term{free-group}{自由群}。
\end{SGAresult}
\begin{proof}
对相交得到 $X$ 的\Term{hypersurface}{超曲面}的数目作归纳，应用3.6，并注意对\Term{dimension}{维数} $\geq3$ 的\Term{complete-intersection}{完全交} $X$，当 $i=1,2$ 及任意 $n$ 时，都有 $H^i(X,\mathcal O_X(n))=0$。
\end{proof}
\textbf{注记 3.8。}当 $X$ 是\Term{nonsingular}{非奇异}\Term{hypersurface}{超曲面}时，3.7归功于 Andreotti。结果3.7（当 $X$ \Term{nonsingular}{非奇异}时）也\Term{representable}{可表}达为：$X$ 的\Term{homogeneous-coordinate-ring}{齐次坐标环}是\Term{factorial}{阶乘}环；在这种形式下，它包含在 XI 3.13(ii) 中。还应指出，Serre 曾在\Term{nonsingular}{非奇异}情形下通过超越方法证明3.7：他用\Term{specialization}{特化}论证归结到\Term{characteristic}{特征} $0$ 的情形，此时可以运用经典形式的 Lefschetz 定理。当然，这里给出的纯\Term{algebra}{代数}证明使我们能够去掉 Lefschetz 定理陈述中的\Term{nonsingular}{非奇异}性假设，这就促使我们也在经典情形中重新考虑这个定理。参见下一讲，其中提出了这一方向的猜想。

在推论3.5和3.7中，我们取基底为\Term{field}{域}，而关键定理2.4和3.4在任意基底上都\Term{effective}{有效}。要把推论3.5和3.6推广到一般的 $S$，就须给出可用的判据，以保证（在 $S$ 上\Term{flat}{平坦}的）$X$ 的一个点具有\Term{pure}{纯}的、分别地\Term{parafactorial}{副阶乘}的\Term{local-ring}{局部环}。这是下一节的主题。

\SourcePage{130}
\SGAsection{4}{形式完备化与法向平坦性}
\begin{SGAresult}{定理}{4.1}
设 $X$ 是\Term{local-scheme}{局部}\Term{noetherian}{诺特}\Term{prescheme}{预概形}，且\Term{local-scheme}{局部}可\Term{immersion}{浸入}\Term{regular}{正则}\Term{scheme}{概形}；$Y$ 是 $X$ 的\Term{closed-subset}{闭子集}，$U=X-Y$；$X_0$ 是由\Term{ideal}{理想} $\mathcal J$ 定义的 $X$ 的闭子\Term{prescheme}{预概形}，$\widehat X$ 是 $X$ 沿 $X_0$ 的\Term{formal-completion}{形式完备化}；$U_0$ 是 $X_0$ 在 $U$ 上的迹，$\widehat U$ 是 $U$ 沿 $U_0$ 的\Term{formal-completion}{形式完备化}，$i:U\to X$ 和 $\widehat i:\widehat U\to\widehat X$ 是\Term{canonical}{典范}\Term{immersion}{浸入}，$n$ 是整数。假设：
\begin{enumerate}
\item[a)] $X$ 在 $Y\cap X_0$ 的各点处沿 $X_0$ \Term{normally-flat}{法向平坦}，即在这些点处，$X_0$ 上的\Term{module}{模} $\mathcal J^n/\mathcal J^{n+1}$ \Term{flat}{平坦}，也就是\Term{locally-free}{局部自由}。
\item[b)] 对每个 $x\in Y\cap X_0$，有 $\operatorname{prof}\mathcal O_{X_0,x}\geq n+2$。
\end{enumerate}
在这些条件下，有：

(1) 设 $F$ 是 $U$ 上的\Term{coherent}{凝聚}\Term{module}{模}，假设：
\begin{enumerate}
\item[c)] 对每个 $x\in Y-Y\cap X_0$，有 $\operatorname{prof}\mathcal O_{X,x}\geq n+2$。
\item[d)] $F$ 在 $U_0$ 的各点处\Term{free}{自由}，且在 $U$ 中它不\Term{free}{自由}的所有点处\Term{depth}{深度} $\geq n+1$。
\end{enumerate}
则 $\bigoplus_{m\geq0}\mathcal J^m$ 上的\Term{graded-module}{分次模}
\[
 \bigoplus_{m\geq0}R^pi_*(\mathcal J^mF)
\]
当 $p\leq n$ 时\Term{finite-type}{有限型}。

(2) 设 $\mathfrak F$ 是 $\widehat U$ 上的\Term{coherent}{凝聚}\Term{module}{模}，则 $\bigoplus_{m\geq0}\mathcal J^m/\mathcal J^{m+1}=\operatorname{gr}_{\mathcal J}(\mathcal O_X)$ 上的\Term{graded-module}{分次模}
\[
 \bigoplus_{m\geq0}R^pi_*(\mathcal J^m\mathfrak F/\mathcal J^{m+1}\mathfrak F)
\]
当 $p\leq n$ 时\Term{finite-type}{有限型}。
\end{SGAresult}
\begin{proof}
1) 令 $X'=\operatorname{Spec}(\bigoplus_{m\geq0}\mathcal J^m)$；\Term{base-change}{基变换} $f:X'\to X$ 定义 $U'=X'-Y'$、$X'_0$、$U'_0=X'_0\cap U'$，以及\Term{immersion}{浸入} $i':U'\to X'$、$i'_0:U'_0\to X'_0$。于是有\Term{cartesian-square}{笛卡尔方块}
\[
\begin{tikzcd}
 X'\arrow[d,"f"'] & U'\arrow[l,"i'"']\arrow[d,"g"]\\
 X & U\arrow[l,"i"']
\end{tikzcd}
\]
并且
\[
 \bigoplus_{m\geq0}R^pi_*(\mathcal J^mF)
 =R^pi_*\left(\bigoplus_{m\geq0}\mathcal J^mF\right)
 =R^pi_*(g_*(F')),\tag{19}
\]
其中 $F'=g^*F$。确实有\Term{canonical}{典范}\Term{isomorphism}{同构}
\[
 g_*(F')\xrightarrow{\sim}\bigoplus_{m\geq0}\mathcal J^mF,
\]
因为在 $U_0$ 的各点处，由 d)，$F$ \Term{free}{自由}，所以成立；在 $U-U_0$ 的各点处，由于 $\mathcal J^m=\mathcal O_U$，也成立（因此在两种情形下，$\mathcal J^m\otimes_{\mathcal O_U}F\to\mathcal J^mF$ 都是\Term{isomorphism}{同构}）。
\SourcePage{131}
另一方面，$f$ 从而 $g$ \Term{affine}{仿射}，因此
\[
 R^pi_*(g_*(F'))=R^p(ig)_*(F')=R^p(fi')_*(F')=f_*(R^pi'_*(F')).\tag{20}
\]
比较(19)、(20)可知，断言1)\Term{equivalence}{等价}于：对每个 $p\leq n$，$R^pi'_*(F')$ 是 $X'$ 上的\Term{finite-type}{有限型}\Term{module}{模}，即\Term{coherent}{凝聚}\Term{module}{模}。由于 $X$ \Term{local-scheme}{局部}可\Term{immersion}{浸入}\Term{regular}{正则}\Term{scheme}{概形}，$X'$ 在 $X$ 上\Term{finite-type}{有限型}，所以 $X'$ 也如此；可以对 $F'$ 的一个\Term{coherent}{凝聚}延拓 $F''$ 应用\Term{coherent}{凝聚}性判据 VIII 2.3：我们要表达 $\mathcal H^p_{Y'}(F'')$ 当 $p\leq n+1$ 时\Term{coherent}{凝聚}，这也\Term{equivalence}{等价}于说，对每个满足
\[
 \operatorname{codim}(\overline{\{x'\}}\cap Y',\overline{\{x'\}})=1\tag{20 bis}
\]
的 $x'\in U'$，都有
\[
 \operatorname{prof}F'_{x'}\geq n+1.\tag{21}
\]
这个条件在 $F'$ 不\Term{free}{自由}的点 $x'$ 处成立，因为由 d)，这样的 $x'$ 不属于 $U'_0$，所以 $g$ 在该点处是\Term{isomorphism}{同构}；再由 d)，$F$ 在 $g(x')$ 处\Term{depth}{深度} $\geq n+1$，因此 $F'$ 在 $x'$ 处\Term{depth}{深度} $\geq n+1$。所以只须在满足(21)且 $F'$ \Term{free}{自由}的 $x'\in U'$ 处验证条件(21)。为此，只须证明在这些点处有
\[
 \operatorname{prof}\mathcal O_{X',x'}\geq n+1.\tag{21 bis}
\]
当然，只要在 $U'$ 中所有满足(20 bis)的点 $x$ 处建立这个关系即可。而再次由第 VIII 讲的判据2.3，这\Term{equivalence}{等价}于断言\Term{module}{模}
\[
 \mathcal H^p_{Y'}(\mathcal O_{X'})\qquad(p\leq n+1)
\]
\Term{coherent}{凝聚}。事实上，我们将证明它们甚至为零；根据第 III 讲，\Term{equivalence}{等价}地说，有
\[
 \operatorname{prof}\mathcal O_{X',x'}\geq n+2\qquad\text{对每个 }x'\in Y'.\tag{22}
\]
为此分两种情形。若 $x'\notin X'_0$，则 $f$ 在 $x'$ 处为\Term{immersion}{浸入}，须验证 $F$ 在像点 $x=f(x')$ 处\Term{depth}{深度} $\geq n+2$，这正是条件 c)。若 $x'\in X'_0$，即 $x=f(x')\in X_0$，从而 $x\in Y\cap X_0$，则借助下面的引理，应用条件 a)、b)。
\end{proof}
\begin{SGAresult}{引理}{4.2}
设 $X$ 是\Term{local-scheme}{局部}\Term{noetherian}{诺特}\Term{prescheme}{预概形}，$X_0$ 是由\Term{ideal}{理想} $\mathcal J$ 定义的 $X$ 的闭子\Term{prescheme}{预概形}，
\[
 X'=\operatorname{Spec}\left(\bigoplus_{m\geq0}\mathcal J^m\right),\qquad
 X'_0=\operatorname{Spec}\left(\bigoplus_{m\geq0}\mathcal J^m/\mathcal J^{m+1}\right)=X'\times_XX_0,
\]
$x$ 是 $X_0$ 中 $X$ 沿 $X_0$ \Term{normally-flat}{法向平坦}的点，即
\[
 \operatorname{gr}_{\mathcal J}(\mathcal O_X)=\bigoplus_{m\geq0}\mathcal J^m/\mathcal J^{m+1}
\]
在该点处作为 $X_0$ 上的\Term{module}{模}\Term{flat}{平坦}。则对 $\mathcal O_{X,x}$ 中任意一列元素 $f_i$（$1\leq i\leq m$），若它们在 $\mathcal O_{X_0,x}$ 中的像构成一条 $\mathcal O_{X_0,x}$-\Term{regular-sequence}{正则序列}，则对 $x$ 上方任意 $x'\in X'$，$f_i$ 在 $\mathcal O_{X',x'}$（分别地，$\mathcal O_{X'_0,x'}$）中的像也构成 $\mathcal O_{X',x'}$-\Term{regular-sequence}{正则序列}（分别地，$\mathcal O_{X'_0,x'}$-\Term{regular-sequence}{正则序列}）。特别地，
\[
 \operatorname{prof}\mathcal O_{X',x'}\geq\operatorname{prof}\mathcal O_{X_0,x},\qquad
 \operatorname{prof}\mathcal O_{X'_0,x'}\geq\operatorname{prof}\mathcal O_{X_0,x}.\tag{23}
\]
\end{SGAresult}
\SourcePage{132}
\begin{proof}
可以假设 $X$ 是以 $x$ 为\Term{closed-point}{闭点}的\Term{local-scheme}{局部}\Term{scheme}{概形}，从而是环为 $A=\mathcal O_{X,x}$ 的\Term{affine}{仿射}\Term{scheme}{概形}，$\mathcal J$ 由\Term{ideal}{理想} $J$ 定义。只须证明，对 $A$ 中每一列元素 $f_i$（$1\leq i\leq m$），若它们在 $A/J$ 中的像构成 $A/J$-\Term{regular-sequence}{正则序列}，则 $f_i$ 同时构成 $\bigoplus_{m\geq0}J^m$-\Term{regular-sequence}{正则序列}和 $\bigoplus_{m\geq0}J^m/J^{m+1}$-\Term{regular-sequence}{正则序列}，即对每个 $m$，构成 $J^m$-\Term{regular-sequence}{正则序列}和 $J^m/J^{m+1}$-\Term{regular-sequence}{正则序列}。第二个断言显然，因为 $J^m/J^{m+1}$ 是\Term{free}{自由} $A/J$-\Term{module}{模}。考察 $J^m$ 的 $J$-进\Term{filtration}{滤过}，并注意 $f_i$ 在 $J^m$ 关于此\Term{filtration}{滤过}的\Term{associated-graded}{相伴分次模}上构成\Term{regular-sequence}{正则序列}，即得第一个断言。
\end{proof}
这证明了4.2，因而证明了4.1的1)。
\begin{proof}[定理4.1的证明（续）]
证明4.1的2)。为此使用\Term{cartesian-square}{笛卡尔方块}
\[
\begin{tikzcd}
 X'_0\arrow[d,"f_0"'] & U'_0\arrow[l,"i'_0"']\arrow[d,"g_0"]\\
 X_0 & U_0\arrow[l,"i_0"']
\end{tikzcd}
\]
并如1)证明的开头那样进行，得到
\[
 \bigoplus_{m\geq0}R^pi_*(\mathcal J^m\mathfrak F/\mathcal J^{m+1}\mathfrak F)
 \simeq f_{0*}(R^pi'_{0*}(\mathfrak F'_0)),\tag{24}
\]
其中 $\mathfrak F_0=\mathfrak F/\mathcal J\mathfrak F$，$\mathfrak F'_0=g_0^*(\mathfrak F_0)$（这里用到了 $\mathfrak F$ \Term{locally-free}{局部自由}）。所以2)的结论\Term{equivalence}{等价}于说，对 $p\leq n$，$R^pi'_{0*}(\mathfrak F'_0)$ 是\Term{coherent}{凝聚}\Term{module}{模}。

这里同样由于 $\mathfrak F'_0$ \Term{locally-free}{局部自由}，判据 VIII 2.3 把问题归结为：以 $\mathcal O_{U'_0}$ 代替 $\mathfrak F'_0$ 时结论成立；即须证明\Term{module}{模}
\[
 \mathcal H^p_{Y'_0}(\mathcal O_{X'_0})\quad(p\leq n+1),\qquad
 Y'_0=Y'\cap X'_0=X'_0-U'_0,
\]
\Term{coherent}{凝聚}。仍然证明它们实际上为零，即
\[
 \operatorname{prof}\mathcal O_{X'_0,x'}\geq n+2\qquad\text{对每个 }x'\in Y'_0.\tag{25}
\]
这确实由条件 a)、b) 及4.2得到。定理4.1得证。
\end{proof}
\textbf{注记 4.3。}由\Term{descent}{下降}立刻可知，假设“$X$ \Term{local-scheme}{局部}可\Term{immersion}{浸入}\Term{regular}{正则}\Term{scheme}{概形}”可以换成下面较弱的假设：存在\Term{faithfully-flat}{忠实平坦}且\Term{quasi-compact}{拟紧}的\Term{morphism}{态射} $\widetilde X\to X$，使 $\widetilde X$ \Term{local-scheme}{局部}可\Term{immersion}{浸入}\Term{regular}{正则}\Term{scheme}{概形}。

定理4.1使我们能够应用第 IX 讲的结果（\Term{comparison-theorem}{比较定理}和\Term{existence-theorem}{存在定理}）。我们特别关心下述推论。
\begin{SGAresult}{推论}{4.4}
假设定理4.1的条件 a)、b)、c) 在 $n=1$ 时成立，且 $X=\operatorname{Spec}(A)$，其中 $A$ 关于 $J$-进\Term{topology}{拓扑}\Term{separated}{分离}且\Term{complete}{完备}。则：
\SourcePage{133}
\begin{enumerate}
\item[(1)] \Term{functor}{函子} $F\mapsto\widehat F$ 从 $U$ 上\Term{locally-free}{局部自由}\Term{coherent}{凝聚}\Term{module}{模}的\Term{category}{范畴}到 $\widehat U$ 上\Term{locally-free}{局部自由}\Term{coherent}{凝聚}\Term{module}{模}的\Term{category}{范畴}\Term{fully-faithful}{完全忠实}。
\item[(2)] 对 $\widehat U$ 上每个\Term{locally-free}{局部自由}\Term{coherent}{凝聚}\Term{module}{模} $\mathfrak F$，存在 $U$ 上的\Term{coherent}{凝聚}\Term{module}{模} $F$ 和\Term{isomorphism}{同构} $\widehat F\xrightarrow{\sim}\mathfrak F$。
\end{enumerate}
特别地，若对 $U$ 中每个其在 $U$ 中的\Term{closure}{闭包}不与 $U_0$ 相交的点 $x$，即满足 $\overline{\{x\}}\cap X_0\subset Y$ 的点 $x$，都有 $\operatorname{prof}\mathcal O_{U,x}\geq2$，则偶对 $(U,U_0)$ 满足第 X 讲的\Term{effective-lefschetz-condition}{有效 Lefschetz 条件} $(\mathrm{Leff})$。
\end{SGAresult}
（最后一条断言同 X 2.1 中一样证明。）

4.4的一个特殊情形是：
\begin{SGAresult}{推论}{4.5}
设 $A$ 是\Term{noetherian}{诺特}环，$J$ 是 $A$ 中包含于\Term{radical}{根基}的\Term{ideal}{理想}，$A_0=A/J$。假设：
\begin{enumerate}
\item[(i)] $\operatorname{prof}A_0\geq3$。
\item[(ii)] $\operatorname{gr}_J(A)$ 是\Term{free}{自由} $A_0$-\Term{module}{模}。
\item[(iii)] $A$ 关于 $J$-进\Term{topology}{拓扑}\Term{complete}{完备}。
\end{enumerate}
设 $X=\operatorname{Spec}(A)$，$X_0=\operatorname{Spec}(A_0)=V(J)$，$a$ 是 $X$ 的\Term{closed-point}{闭点}，$U=X-\{a\}$，$U_0=X_0-\{a\}$，$\widehat U$ 是 $U$ 沿 $U_0$ 的\Term{formal-completion}{形式完备化}。则\Term{functor}{函子} $F\mapsto\widehat F$ 从 $U$ 上\Term{locally-free}{局部自由}\Term{coherent}{凝聚}\Term{module}{模}的\Term{category}{范畴}到 $\widehat U$ 上\Term{locally-free}{局部自由}\Term{coherent}{凝聚}\Term{module}{模}的\Term{category}{范畴}\Term{fully-faithful}{完全忠实}。此外，对 $\widehat U$ 上每个\Term{locally-free}{局部自由}\Term{coherent}{凝聚}\Term{module}{模} $\mathfrak F$，存在 $U$ 上的\Term{coherent}{凝聚}\Term{module}{模}（未必\Term{locally-free}{局部自由}！）$F$ 和\Term{isomorphism}{同构} $\widehat F\simeq\mathfrak F$。
\end{SGAresult}
注意，由4.3，我们无须假设 $A$ 是\Term{regular}{正则}环的商，因为 $A$ 关于 $\operatorname r(A)$-进\Term{topology}{拓扑}的\Term{completion}{完备化}总满足这个条件。

按第 X、XI 讲的方法，由4.5推出：
\begin{SGAresult}{推论}{4.6}
在4.5的条件下，有：
\begin{enumerate}
\item[a)] $U$ 和 $U_0$ \Term{connected}{连通}（III 3.1）。在 $U_0$ 中选定一个\Term{geometric-basepoint}{几何基点}后，\Term{homomorphism}{同态}
\[
 \pi_1(U_0)\longrightarrow\pi_1(U)
\]
是\Term{surjective-map}{满射}。
\item[b)] \Term{homomorphism}{同态} $\operatorname{Pic}(U)\to\operatorname{Pic}(U_0)$ 是\Term{injective-map}{单射}。
\end{enumerate}
\end{SGAresult}
\begin{proof}
证明 b)：考虑到4.5，这归结为验证每个\Term{isomorphism}{同构} $L'_0\xrightarrow{\sim}L_0$ 都可\Term{lifting}{提升}为\Term{isomorphism}{同构} $\widehat L'\xrightarrow{\sim}\widehat L$。为此，逐次\Term{lifting}{提升}为\Term{isomorphism}{同构} $L'_n\xrightarrow{\sim}L_n$；障碍位于
\SourcePage{134}
$H^1(U_0,\mathcal J^n/\mathcal J^{n+1})$ 中，而这些\Term{module}{模}都为零，因为 $J^n/J^{n+1}$ \Term{free}{自由}，且 $\operatorname{prof}A_0\geq3$。
\end{proof}
现在可以证明下述定理。
\begin{SGAresult}{定理}{4.7}
设 $A$ 是\Term{noetherian}{诺特}\Term{local-ring}{局部环}，$J$ 是 $A$ 中包含于其\Term{radical}{根基}的\Term{ideal}{理想}，$A_0=A/J$。假设：
\begin{enumerate}
\item[(i)] $\operatorname{prof}A_0\geq3$。
\item[(ii)] $\operatorname{gr}_J(A)$ 是\Term{free}{自由} $A_0$-\Term{module}{模}。
\end{enumerate}
则若 $A_0$ 是\Term{pure}{纯}的（X 3.1）（分别地，\Term{parafactorial}{副阶乘}的（XI 3.1）），$A$ 也如此。
\end{SGAresult}
\begin{proof}
由\Term{descent}{下降}，可以假设还满足：
\begin{enumerate}
\item[(iii)] $A$ 关于 $J$-进\Term{topology}{拓扑}\Term{complete}{完备}。
\end{enumerate}
事实上，由(i)、(ii)，有 $\operatorname{prof}(A)\geq3$，因而 $\operatorname{prof}(\widehat A)\geq3$，其中 $\widehat A$ 是 $A$ 关于 $J$-进\Term{topology}{拓扑}的\Term{completion}{完备化}；于是应用 X 3.6 和 XI 3.6。因此满足4.5的条件。由于 $\operatorname{prof}(A)\geq3\geq2$，说 $A$ \Term{parafactorial}{副阶乘}就仅仅意味着 $\operatorname{Pic}(U)=0$。由4.6 b)，为此只须 $\operatorname{Pic}(U_0)=0$，即 $A_0$ \Term{parafactorial}{副阶乘}。要证明当 $A_0$ \Term{pure}{纯}时 $A$ 也\Term{pure}{纯}，须证明：若 $V$ 是由 $U$ 上的\Term{algebra}{代数} $\mathcal B$ 定义的 $U$ 的\Term{etale}{平展}\Term{covering}{覆盖}，则 $H^0(U,\mathcal B)$ 是 $A$ 上的\Term{finite}{有限}\Term{etale}{平展}\Term{algebra}{代数}。由于 $A_0$ \Term{pure}{纯}，各 $A_n$ 也\Term{pure}{纯}（它们与 $A_0$ 仅相差\Term{nilpotent}{幂零}元素）；因此对每个 $n$，$B_n=H^0(U,\mathcal B_n)$ 是 $A/J^{n+1}$ 上的\Term{etale}{平展}\Term{algebra}{代数}。这些\Term{algebra}{代数}当然可以\Term{gluing}{粘合}，故 $\varprojlim B_n$ 是 $A$ 上的\Term{etale}{平展}\Term{algebra}{代数}。而由4.5，这个\Term{algebra}{代数}正是 $H^0(U,\mathcal B)$，从而证明了断言。
\end{proof}
\begin{SGAresult}{推论}{4.8}
设 $f:X\to Y$ 是\Term{local-scheme}{局部}\Term{noetherian}{诺特}\Term{prescheme}{预概形}的\Term{flat}{平坦}\Term{morphism}{态射}，$x\in X$，$y=f(x)$。若 $\mathcal O_{X_y,x}$ 是\Term{depth}{深度} $\geq3$ 的纯\Term{local-ring}{局部环}（分别地，\Term{parafactorial}{副阶乘}\Term{local-ring}{局部环}），则 $\mathcal O_{X,x}$ 也如此。
\end{SGAresult}
这就是上一节末为推广推论3.5及其后续结果而许诺的那种结果。因而，使用3.4，得到：
\begin{SGAresult}{推论}{4.9}
设 $f:X\to S$ 是\Term{projective}{射影}且\Term{flat}{平坦}的\Term{morphism}{态射}，$S$ \Term{local-scheme}{局部}\Term{noetherian}{诺特}，$\mathcal O_X(1)$ 是 $X$ 上相对于 $S$ \Term{ample}{丰沛}的\Term{invertible}{可逆}\Term{module}{模}，$t$ 是 $\mathcal O_X(1)$ 的\Term{section}{截面}，且对每个 $s\in S$，它在 $X_s$ 上诱导的\Term{section}{截面} $t_s$ 是 $\mathcal O_{X_s}$-\Term{regular}{正则}的；$X_0$ 是 $t$ 的零点子\Term{scheme}{概形}，$X_m$ 是 $t^{m+1}$ 的零点子\Term{scheme}{概形}。假设对每个 $s\in S$，$X_s$ 在其所有\Term{closed-point}{闭点}处\Term{depth}{深度} $\geq3$。则：
\begin{enumerate}
\item[a)] 若 $X_s-X_{0,s}$（$s\in S$）的\Term{closed-point}{闭点}的\Term{local-ring}{局部环}是\Term{pure}{纯}的，例如为\Term{complete-intersection}{完全交}，则\Term{functor}{函子} $X'\mapsto X'_0=X'\times_XX_0$ 从 $X$ 的\Term{etale}{平展}\Term{covering}{覆盖}\Term{category}{范畴}到 $X_0$ 的\Term{etale}{平展}\Term{covering}{覆盖}\Term{category}{范畴}是\Term{category-equivalence}{范畴等价}；特别地，在 $X_0$ 中选定一个\Term{geometric-basepoint}{几何基点}后，
\SourcePage{135}
\Term{homomorphism}{同态}
\[
 \pi_1(X_0)\longrightarrow\pi_1(X)
\]
是\Term{isomorphism}{同构}。\footnote{编者注：应当提到，后来 Fulton 和 Hansen 在 $S=\operatorname{Spec}(k)$（$k$ 为\Term{algebraically-closed-field}{代数闭域}）的情形得到一个出色的\Term{connectedness}{连通性}结果。设 $g:X\to\mathbf P_k^m\times\mathbf P_k^m$ 满足 $\dim g(X)>m$，则\Term{diagonal}{对角线}的\Term{inverse-image}{逆像}\Term{connected}{连通}。这使人可以推广 $f$ 为 $\mathbf P_k^m$ 在 $S=\operatorname{Spec}(k)$ 上的\Term{structural-morphism}{结构态射}时的推论4.9：确切地说，$\mathbf P_k^m$ 中\Term{dimension}{维数} $>m/2$ 的\Term{irreducible}{不可约}子\Term{variety}{簇} $X$ 的\Term{fundamental-group}{基本群}\Term{trivial-bundle}{平凡}！参见 Fulton W. \& Hansen J.，« A connectedness theorem for projective varieties, with applications to intersections and singularities of mappings »，\emph{Ann. of Math.} (2) \textbf{110} (1979)，no.~1，159--166页。推广到 Grassmann \Term{variety}{簇}或阿贝尔\Term{variety}{簇}的结果，见 Debarre O.，« Théorèmes de connexité pour les produits d'espaces projectifs et les grassmanniennes »，\emph{Amer. J. Math.} \textbf{118} (1996)，no.~6，1347--1367页；以及 « Théorèmes de connexité et variétés abéliennes »，\emph{Amer. J. Math.} \textbf{117} (1995)，no.~3，787--805页。上述 $X$ 的\Term{fundamental-group}{基本群}\Term{trivial-bundle}{平凡}的结果也由 Faltings 独立得到；他还证明 $\operatorname{Pic}(X)$ 没有阶与 $k$ 的\Term{characteristic}{特征}互素的挠元，其方法是\Term{formal-vector-bundle}{形式向量丛}的\Term{algebraization}{代数化}，更接近 Grothendieck 的技术，参见 Faltings G.，« Algebraization of some formal vector bundles »，\emph{Ann. of Math.} (2) \textbf{110} (1979)，no.~3，501--514页。}
\item[b)] 若 $X_s-X_{0,s}$（$s\in S$）的\Term{closed-point}{闭点}的\Term{local-ring}{局部环}\Term{parafactorial}{副阶乘}，例如\Term{regular}{正则}，或为\Term{dimension}{维数} $\geq4$ 的\Term{complete-intersection}{完全交}，则对每个满足 $n>m$ 且 $i=1,2$ 时 $R^if_{0*}(\mathcal O_{X_0}(-n))=0$ 的整数 $m$，映射 $\operatorname{Pic}(X)\to\operatorname{Pic}(X_m)$ 是\Term{bijective-map}{双射}。
\end{enumerate}
此外，若 $S$ \Term{noetherian}{诺特}，且各 $X_{0,s}$ 在其\Term{closed-point}{闭点}处\Term{depth}{深度} $\geq3$，则存在这样的 $m$（参见1.5）。
\end{SGAresult}
\textbf{注记 4.10。}在4.9 b)最后一条断言的条件下，已在1.5中看到，存在 $m$，使 $n>m$ 甚至蕴含 $i=1,2$（甚至 $i\leq2$）时 $H^i(X_0,\mathcal O_{X_0}(-n))=0$。这个条件强于 $i=1,2$ 时 $R^if_*(\mathcal O_{X_0}(-n))=0$，而且有在\Term{base-change}{基变换}下\Term{stable}{稳定}的优点。4.9中所作的\Term{depth}{深度}假设也如此，“各 $X_s$ \Term{local-scheme}{局部}为\Term{complete-intersection}{完全交}”一类假设也如此。因此，在这些条件下，4.9 b)还蕴含 $\mathrm{Sch}/S$ 上的\Term{functor}{函子}\Term{morphism}{态射}
\[
 \operatorname{Pic}_{X/S}\longrightarrow\operatorname{Pic}_{X_n/S}
\]
是\Term{isomorphism}{同构}，因而相对 \Term{picard-scheme}{Picard 概形}存在时，其间的\Term{morphism}{态射}
\[
 \mathbf{Pic}_{X/S}\longrightarrow\mathbf{Pic}_{X_n/S}
\]
也是\Term{isomorphism}{同构}。即使 $S$ 是\Term{algebraically-closed-field}{代数闭域}的谱，这个命题也比仅仅说 $\operatorname{Pic}(X)\to\operatorname{Pic}(X_n)$ 是\Term{bijective-map}{双射}精确得多。

\SourcePage{136}
可以问，在前面的结论中，能否总取 $n=0$（因此假设各 $X_{0,s}$ 在其\Term{closed-point}{闭点}处\Term{depth}{深度} $\geq3$）。当 $X_0$ 在 $S$ 上\Term{smooth}{光滑}，且 $S$ 的\Term{residue-characteristic}{剩余特征}都为零时，确实如此。这来自 Kodaira 的“消失定理”（由超越方法、利用 Kähler 度量证明）：它蕴含，对\Term{characteristic}{特征}为零的\Term{field}{域} $k$ 上每个\Term{dimension}{维数}为 $n$ 的\Term{smooth}{光滑}\Term{connected}{连通}\Term{projective}{射影}\Term{scheme}{概形}，以及 $X$ 上每个\Term{ample}{丰沛}\Term{invertible}{可逆}\Term{module}{模} $L$，当 $i\neq n$ 时都有 $H^i(X,L^{-1})=0$。目前尚不清楚\footnote{编者注：正如 Raynaud 指出的，Deligne 和 Illusie 关于 de Rham \Term{complex}{复形}\Term{resolution}{分解}的结果容易推出：若假设\Term{characteristic}{特征} $p>0$ 的 $k$ 上的\Term{smooth}{光滑}\Term{projective}{射影} $X$ 可以\Term{lifting}{提升}为 $W_2(k)$ 上的\Term{flat}{平坦}\Term{scheme}{概形}，则对 $X$ 上\Term{ample}{丰沛}的 $L$，当 $i<\inf(p,\dim(X))$ 时，群 $H^i(X,L^{-1})$ 为零。参见 Deligne P. \& Illusie L.，« Relèvements modulo $p^2$ et décomposition du complexe de de Rham »，\emph{Invent. Math.} \textbf{89} (1987)，no.~2，247--270页。这给出\Term{characteristic}{特征}零时\Term{projective}{射影}\Term{variety}{簇}的 Kodaira 结果的纯\Term{algebra}{代数}证明。若 $X$ 不可\Term{lifting}{提升}，众所周知“消失定理”不成立；参见 Raynaud M.，« Contre-exemple au ”vanishing theorem” en caractéristique $p>0$ »，载 \emph{C.P. Ramanujam---a tribute}，Tata Inst. Fund. Res. Studies in Math.，第8卷，Springer，Berlin--New York，1978，273--278页中的例子；另见 Haboush W. \& Lauritzen N.，« Varieties of unseparated flags »，载 \emph{Linear algebraic groups and their representations (Los Angeles, CA, 1992)}，Contemp. Math.，第153卷，American Mathematical Society，Providence，RI，1993，35--57页中的精彩例子，其简化见 Lauritzen N. \& Rao A.P.，« Elementary counterexamples to Kodaira vanishing in prime characteristic »，\emph{Proc. Indian Acad. Sci. Math. Sci.} \textbf{107} (1997)，no.~1，21--25页。反之，我不知道正\Term{characteristic}{特征}下 $n>1$ 时\Term{arrow}{箭头} $\operatorname{Pic}(X_{n+1})\to\operatorname{Pic}(X_n)$ 非\Term{surjective-map}{满射}的例子，其中 $X_n$ 是上述\Term{smooth}{光滑}\Term{projective}{射影} $X$ 的加厚\Term{hyperplane-section}{超平面截面}。}，这个定理是否能换成\Term{characteristic}{特征} $p>0$ 中的推广，以及\Term{smoothness}{光滑性}假设是否能换成更一般性质的假设（关于\Term{depth}{深度}，或“\Term{complete-intersection}{完全交}”类型的假设等）。

\SGAsection{5}{非固有态射的普遍有限性条件}
为便于参照，回顾如下命题。
\begin{SGAresult}{命题}{5.1}
设 $f:X\to S$ 是\Term{prescheme}{预概形}的\Term{proper}{固有}\Term{morphism}{态射}，$S$ \Term{local-scheme}{局部}\Term{noetherian}{诺特}，$U$ 是 $X$ 的\Term{open-subset}{开子集}，$g:U\to X$ 是\Term{canonical}{典范}\Term{immersion}{浸入}，$h=fg:U\to S$，$F$ 是 $U$ 上的\Term{module}{模}。假设 $i\leq n$ 时\Term{module}{模} $R^ig_*(F)$ \Term{coherent}{凝聚}（这是 $X$ 上\Term{local-scheme}{局部}的假设，实际可用判据 VIII 2.3 验证）。则 $i\leq n$ 时 $R^ih_*(F)$ \Term{coherent}{凝聚}。
\end{SGAresult}
这立即来自 Leray \Term{spectral-sequence}{谱序列}
\[
 E_2^{p,q}=R^pf_*(R^qg_*(F))\Longrightarrow R^*h_*(F),\tag{26}
\]
以及 $X$ 上\Term{coherent}{凝聚}\Term{module}{模}在 $f$ 下的高阶\Term{direct-image}{正像}\Term{coherent}{凝聚}这一事实（EGA III 3.2.1）。

\SourcePage{137}
\begin{SGAresult}{命题}{5.2}
设 $S$ 是\Term{local-scheme}{局部}\Term{noetherian}{诺特}\Term{prescheme}{预概形}，$\mathscr S$ 是 $S$ 上\Term{finite-type}{有限型}、由 $\mathscr S_1$ 生成的\Term{quasi-coherent}{拟凝聚}\Term{graded-algebra}{分次代数}，$X$ 是 $\operatorname{Proj}(\mathscr S)$ 的子\Term{prescheme}{预概形}，$\mathcal O_X(1)$ 是 $\operatorname{Proj}(\mathscr S(1))$ 诱导的 $X$ 上相对于 $S$ \Term{very-ample}{甚丰沛}的\Term{invertible}{可逆}\Term{module}{模}，$U$ 是 $X$ 的\Term{open-subset}{开子集}，$g:U\to X$ 是\Term{canonical}{典范}\Term{immersion}{浸入}，$h=fg:U\to S$，$F$ 是 $U$ 上的\Term{quasi-coherent}{拟凝聚}\Term{module}{模}，从而得到扭转\Term{module}{模} $F(m)=F\otimes\mathcal O_X(m)$（$m\in\mathbf Z$），$n$ 是整数，$m_0$ 是整数。下列条件\Term{equivalence}{等价}：
\begin{enumerate}
\item[(i)] 当 $i\leq n$ 时，$R^ig_*(F)$ \Term{coherent}{凝聚}。
\item[(ii)] 当 $i\leq n$ 时，$\bigoplus_{m>m_0}R^ih_*(F(m))$ 是\Term{finite-type}{有限型} $\mathscr S$-\Term{module}{模}。
\end{enumerate}
\end{SGAresult}
\begin{proof}
在上述\Term{spectral-sequence}{谱序列}中以 $F(m)$ 代替 $F$，得到\Term{graded}{分次} $\mathscr S$-\Term{module}{模}的\Term{spectral-sequence}{谱序列}
\[
 E_2^{p,q}=\bigoplus_{m\geq m_0}R^pf_*(R^qg_*(F(m)))
 \Longrightarrow\bigoplus_{m\geq m_0}R^*h_*(F(m)).
\]
由于
\[
 R^qg_*(F(m))\simeq R^qg_*(F)(m),
\]
可见若 $R^ig_*(F)$ \Term{coherent}{凝聚}，则由下面引理5.3的 a) 部分，$q\leq n$ 时 $E_2^{p,q}$ 在 $\mathscr S$ 上\Term{finite-type}{有限型}，因而终项在次数 $i\leq n$ 时在 $\mathscr S$ 上\Term{finite-type}{有限型}。这证明了(i)$\Rightarrow$(ii)。此外，在\Term{graded}{分次} $\mathscr S$-\Term{module}{模}的\Term{abelian-category}{阿贝尔范畴}对由\Term{quasi-coherent}{拟凝聚}\Term{finite-type}{有限型}\Term{module}{模}组成的\Term{thick-subcategory}{厚子范畴} $\mathcal C$ 取商所得的\Term{category}{范畴}中论证，上述\Term{spectral-sequence}{谱序列}给出
\[
 \bigoplus_{m\geq m_0}R^{n+1}h_*(F(m))
 \simeq\bigoplus_{m\geq m_0}f_*(R^{n+1}g_*(F)(m))\pmod{\mathcal C}.
\]
由引理5.3的 b) 部分，这表明若左边是\Term{finite-type}{有限型} $\mathscr S$-\Term{module}{模}，则 $R^{n+1}g_*(F)$ \Term{coherent}{凝聚}。对 $n$ 归纳，便证明了(ii)$\Rightarrow$(i)。还须证明下述引理。
\end{proof}
\begin{SGAresult}{引理}{5.3}
设 $S,\mathscr S,X,f$ 如5.2，$G$ 是 $X$ 上的\Term{quasi-coherent}{拟凝聚}\Term{module}{模}，$m_0$ 是整数。则：
\begin{enumerate}
\item[a)] 若 $G$ \Term{coherent}{凝聚}，则对每个整数 $i$，$\mathscr S$ 上的\Term{graded-module}{分次模}
\[
 \bigoplus_{m\geq m_0}R^if_*(G(m))
\]
\Term{finite-type}{有限型}。
\item[b)] 反之，假设 $\mathscr S$ 上的\Term{module}{模} $\bigoplus_{m\geq m_0}R^if_*(G(m))$ \Term{finite-type}{有限型}，则 $G$ \Term{coherent}{凝聚}。
\end{enumerate}
\end{SGAresult}
\begin{proof}[引理5.3的证明]
对于 a)，$i=0$ 的情形见 EGA III 2.3.2；$i>0$ 的情形见 EGA III 2.2.1(i)(ii)，它说 $R^if_*G(m)$ \Term{coherent}{凝聚}，且当 $m$ 充分大时为零（这里可假设 $S$ \Term{noetherian}{诺特}）。对于 b)，注意 $G$ \Term{isomorphism}{同构}于
\[
 \operatorname{Proj}\left(\bigoplus_{m\geq m_0}f_*(G(m))\right)
\]
（EGA II 3.4.4和3.4.2），由该处3.4.4可知，若 $\bigoplus_{m\geq m_0}f_*(G(m))$ 是\Term{finite-type}{有限型} $\mathscr S$-\Term{module}{模}，则 $G$ \Term{coherent}{凝聚}。
\end{proof}

\SourcePage{138}
\begin{SGAresult}{推论}{5.4}
（$S$ \Term{noetherian}{诺特}。）假设 $i\leq n$ 时 $R^ig_*(F)$ \Term{coherent}{凝聚}，则当 $i\leq n+1$ 且 $m$ 充分大时，有\Term{canonical}{典范}\Term{isomorphism}{同构}
\[
 R^ih_*(F(m))\simeq f_*(R^ig_*(G)(m)).
\]
\end{SGAresult}
事实上，由 EGA III 2.2.1(ii)，用于 $F(m)$ 的\Term{spectral-sequence}{谱序列}(26)在次数 $\leq n$ 时\Term{degeneration}{退化}，立刻得到结果（这也重新给出了5.2中(ii)$\Rightarrow$(i)）。
\begin{SGAresult}{推论}{5.5}
在5.2的预备条件下，设 $S$ \Term{noetherian}{诺特}，则下列条件\Term{equivalence}{等价}：
\begin{enumerate}
\item[(i)] $\bigoplus_{m\geq m_0}h_*(F(m))$ 在 $S$ 上\Term{finite-type}{有限型}，并且当 $0<i\leq n$ 且 $m$ 充分大时，$R^ih_*(F(m))=0$。
\item[(ii)] $g_*(F)$ \Term{coherent}{凝聚}，且当 $0<i\leq n$ 时，$R^ig_*(F)=0$。
\item[(ii bis)] $g_*(F)$ \Term{coherent}{凝聚}，且 $\operatorname{prof}_Yg_*(F)>n+1$。
\end{enumerate}
\end{SGAresult}
(ii)与(ii bis)的\Term{equivalence}{等价}性包含在 III 3.3 中。此外，由5.2，条件(i)、(ii)都蕴含各 $R^ig_*(F)$（$i\leq n$）\Term{coherent}{凝聚}。于是(i)与(ii)的\Term{equivalence}{等价}性由5.4得到，只要考虑到：对 $X$ 上每个\Term{coherent}{凝聚}\Term{module}{模} $G$，$G=0$ 当且仅当 $m$ 充分大时 $f_*(G(m))=0$；例如可由 EGA III 2.2.1(iii)得到这一事实。

\textbf{注记 5.6。}可以如下解释判据5.2和5.5：“同时\Term{finite}{有限}性条件”5.2(ii)\Term{representable}{可表}达为 $F$ 在 $U$ 中邻近 $Y=X-U$ 的点处的\Term{local-scheme}{局部}\Term{regular}{正则}性性质（借助 VIII 2.1，用\Term{depth}{深度}表达）；而“渐近消失条件”5.5(i)则强得多，\Term{representable}{可表}达为 $g_*(F)$ 在 $Y$ 本身各点处的\Term{local-scheme}{局部}\Term{regular}{正则}性条件。为了把\Term{projective}{射影}\Term{morphism}{态射}的 Lefschetz 型定理推广到\Term{quasi-projective}{拟射影}\Term{morphism}{态射}，有必要在 $X$ 上寻找必要充分的\Term{local-scheme}{局部}判据，使 $\mathscr S$-\Term{module}{模}
\[
 \bigoplus_{m\geq0}R^ih_*(F(m))\qquad(i\leq n)
\]
\Term{finite-type}{有限型}。当 $S$ 是\Term{field}{域}的谱（更一般地，若是\Term{artinian-ring}{阿廷环}的谱，想来也如此），并且 $Y=X-U$ \Term{finite}{有限}时，可以证明，以下条件成立是必要且充分的：
\begin{enumerate}
\item[$1^\circ$)] 对 $U$ 的每个\Term{closed-point}{闭点} $x$，$\operatorname{prof}F_x>n$（比较1.4）。
\item[$2^\circ$)] 当 $i\leq n$ 时，$R^ig_*(F)$ \Term{coherent}{凝聚}；或\Term{equivalence}{等价}地，存在 $Y$ 的\Term{open-neighborhood}{开邻域} $V$，使对 $U\cap V$ 的每个\Term{closed-point}{闭点} $x$，都有 $\operatorname{prof}F_x>n+1$。
\end{enumerate}
\begin{SGAresult}{命题}{5.7}
设 $S$ 是\Term{local-scheme}{局部}\Term{noetherian}{诺特}\Term{prescheme}{预概形}，$g:U\to X$ 是 $S$ 上\Term{finite-type}{有限型}\Term{prescheme}{预概形}间的\Term{morphism}{态射}\footnote{事实上，只须 $g$ \Term{quasi-compact}{拟紧}且\Term{quasi-separated}{拟分离}（EGA IV 1.2.1），对 $U,X$ 不加条件。}，$h,f$ 是其\Term{structural-morphism}{结构态射}，$F$ 是 $U$ 上的\Term{quasi-coherent}{拟凝聚}\Term{module}{模}，$n$ 是整数。下列条件\Term{equivalence}{等价}：
\begin{enumerate}
\item[(i)] 对每个\Term{base-change}{基变换} $S'\to S$，其中 $S'$ \Term{noetherian}{诺特}，$X'$ 上的\Term{module}{模} $R^ng'_*(F')$ \Term{coherent}{凝聚}。
\SourcePage{139}
\item[(ii)] 对每个如上的\Term{base-change}{基变换}以及 $S'$ 上每个\Term{coherent}{凝聚}\Term{ideal}{理想} $\mathcal J$，以 $\mathcal I$ 表示 $X'$ 上的\Term{ideal}{理想} $\mathcal J\mathcal O_{X'}$，则 $\bigoplus_{m\geq0}\mathcal I^m$ 上的\Term{graded-module}{分次模}
\[
 \bigoplus_{m\geq0}R^ng'_*(\mathcal I^mF')
\]
\Term{finite-type}{有限型}。
\item[(iii)] 对每个\Term{base-change}{基变换} $S'\to S$ 以及如上的 $\mathcal J$，$\operatorname{gr}_{\mathcal I}(\mathcal O_{X'})=\bigoplus_{m\geq0}\mathcal I^m/\mathcal I^{m+1}$ 上的\Term{graded-module}{分次模}
\[
 \bigoplus_{m\geq0}R^ng'_*(\mathcal I^mF'/\mathcal I^{m+1}F')
\]
\Term{finite-type}{有限型}。
\end{enumerate}
\end{SGAresult}
\begin{proof}
显然(ii)$\Rightarrow$(i)且(iii)$\Rightarrow$(i)，只须在条件(ii)、(iii)中取 $\mathcal I=0$ 即可看出。反向蕴含则把(i)应用于复合\Term{base-change}{基变换} $S''\to S'\to S$ 得到，其中分别取
\[
 S''=\operatorname{Spec}\left(\bigoplus_{m\geq0}\mathcal J^m\right),\qquad
 S''=\operatorname{Spec}\left(\bigoplus_{m\geq0}\mathcal J^m/\mathcal J^{m+1}\right).
\]
\end{proof}
这个命题的意义在于：(ii)形式的条件正是“\Term{algebra}{代数}与形式\Term{comparison-theorem}{比较定理}”中出现的条件，而(iii)形式的条件出现在补充它们的“\Term{existence-theorem}{存在定理}”中，参见第 IX 讲。第一个有意义的情形是 $f:X\to S$ 为恒等\Term{morphism}{态射}；这时所讨论的是对\Term{local-scheme}{局部}\Term{finite-type}{有限型}\Term{morphism}{态射} $h:U\to S$ 以及 $U$ 上相对于 $S$ \Term{flat}{平坦}的\Term{quasi-coherent}{拟凝聚}\Term{module}{模} $F$ 的条件。为得到充分条件，我们将假设 $U$ 通过 $g:U\to X$ 嵌入为 $S$ 上\Term{proper}{固有}的 $X$ 的开子\Term{prescheme}{预概形}。应用5.1，就得到：
\begin{SGAresult}{推论}{5.8}
设 $f:X\to S$ 是\Term{proper}{固有}\Term{morphism}{态射}，$S$ \Term{local-scheme}{局部}\Term{noetherian}{诺特}，$U$ 是 $X$ 的\Term{open-subset}{开集}，$g:U\to X$ 是\Term{canonical}{典范}\Term{immersion}{浸入}，$h=fg:U\to S$，$F$ 是 $X$ 上相对于 $S$ \Term{flat}{平坦}的\Term{quasi-coherent}{拟凝聚}\Term{module}{模}。假设对每个\Term{base-change}{基变换} $S'\to S$，其中 $S'$ \Term{local-scheme}{局部}\Term{noetherian}{诺特}，当 $i\leq n$ 时，$R^ig'_*(F')$ 在 $X'$ 上\Term{coherent}{凝聚}。则有：
\begin{enumerate}
\item[(i)] 对每个\Term{base-change}{基变换} $S'\to S$，其中 $S'$ \Term{local-scheme}{局部}\Term{noetherian}{诺特}，当 $i\leq n$ 时，$R^ih'_*(F')$ 在 $S'$ 上\Term{coherent}{凝聚}。
\item[(ii)] 对每个如上的 $S'\to S$，以及 $S'$ 上每个\Term{coherent}{凝聚}\Term{ideal}{理想} $\mathcal J$，\Term{graded-module}{分次模}
\[
 \bigoplus_{m\geq0}R^ih'_*(\mathcal J^mF')
\]
当 $i\leq n$ 时在 $\bigoplus_{m\geq0}\mathcal J^m$ 上\Term{finite-type}{有限型}。
\item[(iii)] 对每个如上的 $S'\to S$ 和 $\mathcal J$，\Term{graded-module}{分次模}
\[
 \bigoplus_{m\geq0}R^ih'_*(\mathcal J^mF'/\mathcal J^{m+1}F')
\]
当 $i\leq n$ 时在 $\operatorname{gr}_{\mathcal J}(\mathcal O_{S'})=\bigoplus_{m\geq0}\mathcal J^m/\mathcal J^{m+1}$ 上\Term{finite-type}{有限型}。
\end{enumerate}
\SourcePage{140}
此外，在(ii)的条件下，根据\Term{comparison-theorem}{比较定理}1.1，以 $\widehat S'$ 表示 $S'$ 关于 $\mathcal J$ 的\Term{formal-completion}{形式完备化}，以 $\widehat U'$ 表示 $U'$ 关于 $\mathcal J\mathcal O_{U'}$ 的\Term{formal-completion}{形式完备化}，则\Term{canonical}{典范}\Term{homomorphism}{同态}
\[
 \widehat{R^ih'_*(F')}\longrightarrow R^i\widehat h'_* (\widehat F')
 \longrightarrow\varprojlim_k R^ih'_*(F'_k)
\]
当 $i\leq n-1$ 时都是\Term{isomorphism}{同构}。
\end{SGAresult}
\textbf{注记 5.9。}再假设满足5.8的条件，$F$ \Term{coherent}{凝聚}，并考虑如5.9(i)中的\Term{base-change}{基变换} $S'\to S$。再假设 $S'$ \Term{local-scheme}{局部}可\Term{immersion}{浸入}\Term{regular}{正则}\Term{scheme}{概形}；或更一般地，存在\Term{faithfully-flat}{忠实平坦}且\Term{quasi-compact}{拟紧}的\Term{morphism}{态射} $S''\to S$，使 $S''$ \Term{local-scheme}{局部}可\Term{immersion}{浸入}\Term{regular}{正则}\Term{scheme}{概形}；特别当 $S'$ 为\Term{local-scheme}{局部}\Term{scheme}{概形}时，这个条件成立。则把 $F'$ 换成 $U'$ 上的\Term{module}{模} $G'$ 后，5.8(i)、(ii)的结论仍成立，只要 $U'$ 的每个点都有一个\Term{open-neighborhood}{开邻域}，其上 $G'$ \Term{isomorphism}{同构}于形如 $F'{}^n$ 的\Term{module}{模}。事实上，可归结到 $S'$ 本身\Term{local-scheme}{局部}可\Term{immersion}{浸入}\Term{regular}{正则}\Term{scheme}{概形}的情形，于是
\[
 S''=\operatorname{Spec}\left(\bigoplus_{m\geq0}\mathcal J^m\right),\qquad
 X''=X'\times_{S'}S''=X\times_SS''
\]
也如此，因为它们在其上\Term{finite-type}{有限型}。然后对 $G''$ 在\Term{immersion}{浸入} $U''\to X''$ 下次数 $i\leq n$ 的\Term{direct-image}{正像}应用\Term{finite}{有限}性判据 VIII 2.3，注意这些条件由假设对 $F''$ 成立，因而也对 $G''$ 成立，因为它们用\Term{depth}{深度}表达，而 $G''$ \Term{local-scheme}{局部}\Term{isomorphism}{同构}于某个 $F''{}^n$。同样的论证表明：若 $\mathscr G'$ 是 $\widehat U'$（$U'$ 关于\Term{ideal}{理想} $\mathcal J\mathcal O_{U'}$ 的\Term{completion}{完备化}）上的\Term{coherent}{凝聚}\Term{module}{模}，且 $\mathscr G'_0=\mathscr G'/\mathcal J\mathscr G'$ \Term{local-scheme}{局部}形如 $F'_0{}^n$，则把 $F'$ 换成 $\mathscr G'$ 后，(iii)的结论仍成立。利用第 IX 讲的结果，得到：
\begin{SGAresult}{推论}{5.10}
设 $f:X\to S$ 是\Term{proper}{固有}\Term{morphism}{态射}，$S$ \Term{local-scheme}{局部}\Term{noetherian}{诺特}，$U$ 是 $X$ 的\Term{open-subset}{开子集}；假设 $U$ 相对于 $S$ \Term{flat}{平坦}，并且对每个\Term{base-change}{基变换} $S'\to S$，其中 $S'$ \Term{local-scheme}{局部}\Term{noetherian}{诺特}，当 $i=0,1$ 时，$R^ig'_*(\mathcal O_{U'})$ 在 $X'$ 上\Term{coherent}{凝聚}。再假设 $S'$ 形如 $\operatorname{Spec}(A)$，其中 $A$ 是具有\Term{ideal}{理想} $J$ 的\Term{noetherian}{诺特}环，且 $A$ 关于 $J$-进\Term{topology}{拓扑}\Term{separated}{分离}且\Term{complete}{完备}。在这些条件下：
\begin{enumerate}
\item[(i)] \Term{functor}{函子} $F\mapsto\widehat F$ 从 $U'$ 上\Term{locally-free}{局部自由}\Term{module}{模}的\Term{category}{范畴}到 $\widehat U'$ 上\Term{locally-free}{局部自由}\Term{module}{模}的\Term{category}{范畴}\Term{fully-faithful}{完全忠实}。
\item[(ii)] 对 $\widehat U'$ 上每个\Term{locally-free}{局部自由}\Term{module}{模} $\mathfrak F$，存在 $U'$ 上的\Term{coherent}{凝聚}\Term{module}{模} $F$（遗憾的是，未必\Term{locally-free}{局部自由}），以及\Term{isomorphism}{同构} $\widehat F\simeq\mathfrak F$。
\end{enumerate}
\end{SGAresult}
\begin{proof}
由5.9，只须再证明(ii)。根据这个注记和2.1，$\mathfrak F$ 由 $\widehat X'$ 上的\Term{coherent}{凝聚}\Term{module}{模} $\mathfrak G$ 诱导。由\Term{existence-theorem}{存在定理} EGA III 5.1.4，$\mathfrak G$ 形如 $\widehat F$，其中 $F$ 在 $X$ 上\Term{coherent}{凝聚}，从而得到结论。
\end{proof}

\SourcePage{141}
\textbf{注记 5.11。} $1^\circ$）利用5.10、4.7以及一个适当假设——$X'\to S'$ 的几何\Term{fibre}{纤维}的某些\Term{local-ring}{局部环}是\Term{pure}{纯}的，分别地\Term{parafactorial}{副阶乘}的——应该可以得到如下命题：\Term{functor}{函子} $Z'\mapsto\widehat Z'$ 从 $X'$ 的\Term{etale}{平展}\Term{covering}{覆盖}\Term{category}{范畴}到 $\widehat X'$ 的\Term{etale}{平展}\Term{covering}{覆盖}\Term{category}{范畴}（\Term{equivalence}{等价}地，到 $X'_0$ 的\Term{etale}{平展}\Term{covering}{覆盖}\Term{category}{范畴}）是\Term{category-equivalence}{范畴等价}；分别地，\Term{functor}{函子} $L\mapsto\widehat L$ 从 $X'$ 上\Term{invertible}{可逆}\Term{module}{模}的\Term{category}{范畴}到 $\widehat X'$ 上\Term{invertible}{可逆}\Term{module}{模}的\Term{category}{范畴}是\Term{category-equivalence}{范畴等价}。利用 Murre 最近的结果，很可能应当能推出某些非\Term{proper}{固有}\Term{algebraic-scheme}{代数概形}的 \Term{picard-scheme}{Picard 概形}的\Term{existence-theorem}{存在定理}。\footnote{编者注：当然，在\Term{projective}{射影}情形下，应参看 Grothendieck 在 FGA 中的\Term{existence-theorem}{存在定理}；参见 Grothendieck A.，« Technique de descente et théorèmes d'existence en géométrie algébrique. VI. Les schémas de Picard : propriétés générales »，载 \emph{Séminaire Bourbaki}，第7卷，Société mathématique de France，Paris，1995，Exp.~236，221--243页；以及 « Technique de descente et théorèmes d'existence en géométrie algébrique. V. Les schémas de Picard : théorèmes d'existence »，载同卷，1995，Exp.~232，143--161页。那里提出的九个\Term{finite}{有限}性猜想，由 Raynaud 女士和 Kleiman 在 SGA 6 的第 XII、XIII 讲中证明。关于这个主题，一篇出色的初等文本是 Kleiman 的综述：Kleiman S.，« The Picard scheme »，即将发表于 \emph{Contemp. Math.}。这些技术对于相对\Term{smooth}{光滑}曲线的整体广义 Jacobian 的应用，见 Contou-Carrère C.，« La jacobienne généralisée d'une courbe relative ; construction et propriété universelle de factorisation »，\emph{C. R. Acad. Sci. Paris Sér. A-B} \textbf{289} (1979)，no.~3，A203--A206页；以及 « Jacobiennes généralisées globales relatives »，载 \emph{The Grothendieck Festschrift}，Vol.~II，Progr. Math.，第87卷，Birkhäuser，Boston，1990，69--109页。另见同一作者在纯\Term{local-scheme}{局部}框架中对“\Term{local-scheme}{局部}广义 Jacobian”\Term{functor}{函子}的构造与研究：« Jacobienne locale, groupe de bivecteurs de Witt universel, et symbole modéré »，\emph{C. R. Acad. Sci. Paris Sér. I Math.} \textbf{318} (1994)，no.~8，743--746页。另一方面，当\Term{morphism}{态射}\Term{projective}{射影}且\Term{smooth}{光滑}时，\Term{picard-scheme}{Picard 概形}的\Term{connected-component}{连通分支}\Term{proper}{固有}；\Term{singular}{奇异}情形则不再如此。于是自然出现 \Term{picard-scheme}{Picard 概形}的\Term{compactification}{紧化}问题。这个问题已有详细研究，特别见 Altman A.B. \& Kleiman S.，« Compactifying the Picard scheme »，\emph{Adv. in Math.} \textbf{35} (1980)，no.~1，50--112页，以及 « Compactifying the Picard scheme. II »，\emph{Amer. J. Math.} \textbf{101} (1979)，no.~1，10--41页。曲线的情形此前已有研究：D'Souza C.，« Compactification of generalised Jacobians »，\emph{Proc. Indian Acad. Sci. Sect. A Math. Sci.} \textbf{88} (1979)，no.~5，419--457页。甚至已经确切知道一条曲线的\Term{compactification}{紧化} Jacobian 何时\Term{irreducible}{不可约}：Rego C.J.，« The compactified Jacobian »，\emph{Ann. Sci. Éc. Norm. Sup.} (4) \textbf{13} (1980)，no.~2，211--223页；当曲线\Term{geometrically-integral}{几何整}且\Term{local-scheme}{局部}平面时，它就是（通常的）Jacobian 的\Term{closure}{闭包}。族中的\Term{compactification}{紧化} Jacobian 的构造，见 Esteves E.，« Compactifying the relative Jacobian over families of reduced curves »，\emph{Trans. Amer. Math. Soc.} \textbf{353} (2001)，no.~8，3045--3095页。自 SGA 2 初版以来，\Term{proper}{固有}情形的 \Term{picard-scheme}{Picard 概形}存在性结果已有进展；参见 Murre J.P.，« On contravariant functors from the category of pre-schemes over a field into the category of abelian groups (with an application to the Picard functor) »，\emph{Publ. Math. Inst. Hautes Études Sci.} \textbf{23} (1964)，5--43页；尤其是 Artin M.，« Algebraization of formal moduli. I »，载 \emph{Global Analysis (Papers in Honor of K. Kodaira)}，Univ. Tokyo Press，Tokyo，1969，21--71页。另见 Raynaud M.，« Spécialisation du foncteur de Picard »，\emph{Publ. Math. Inst. Hautes Études Sci.} \textbf{38} (1970)，27--76页，讨论\Term{discrete-valuation-ring}{离散赋值环}上\Term{proper}{固有}但未必\Term{flat}{平坦}的\Term{scheme}{概形}。关于更近的结果，特别是 Artin 关于\Term{proper}{固有}且\Term{flat}{平坦}\Term{scheme}{概形}的 \Term{picard-functor}{Picard 函子}的结果，尤其在\Term{dimension}{维数} $0$ 上\Term{cohomology}{上同调}\Term{flat}{平坦}的情形，见 Bosch S.、Lütkebohmert W.、Raynaud M.，\emph{Néron models}，Ergebnisse der Mathematik und ihrer Grenzgebiete (3)，第21卷，Springer-Verlag，Berlin，1990，第 VIII 章及其中参考文献。更近一些，对于\Term{perfect}{完美}\Term{residue-field}{剩余域}的\Term{discrete-valuation-ring}{离散赋值环}的谱 $S$ 上的相对曲线 $f:X\to S$，已经得到非常精细的结果。更确切地，假设 $f$ \Term{proper}{固有}且\Term{flat}{平坦}、$X$ \Term{regular}{正则}，且 $f_*\mathcal O_X=\mathcal O_S$。但不假设 $f$ 在\Term{dimension}{维数} $0$ 上\Term{cohomology}{上同调}\Term{flat}{平坦}，即不假设 $H^1(X,\mathcal O)$ \Term{torsion-free}{无挠}。只要所说的挠部分非零，\Term{picard-scheme}{Picard 概形}就既不能由\Term{scheme}{概形}也不能由\Term{algebraic-space}{代数空间}表示。设 $J$ 是 $f$ 的\Term{generic-fibre}{一般纤维}的 \Term{neron-model}{Néron 模型}；它是 \Term{picard-functor}{Picard 函子} $P$ 的一个商。Raynaud 证明切映射 $H^1(X,\mathcal O)=\operatorname{Lie}(P)\to\operatorname{Lie}(J)$ 的核等于 $H^1$ 的\Term{torsion-subgroup}{挠子群}，且\Term{cokernel}{余核}有相同长度（见 Liu Q.、Lorenzini D.、Raynaud M.，« Néron models, Lie algebras, and reduction of curves of genus one »，\emph{Invent. Math.} \textbf{157} (2004)，455--518页，定理3.1）。这个结果使上述作者能够研究 Birch--Swinnerton-Dyer 猜想与 Artin--Tate 猜想的联系（见该文定理6.6）。关于\Term{local-scheme}{局部} \Term{picard-scheme}{Picard 概形}，见第149页编者注(13)所引 Boutot 的学位论文。}
一般而言，
\SourcePage{142}
利用本讨论班所发展的技术，去掉各种\Term{existence-theorem}{存在定理}中的\Term{purity}{纯性}假设，尤其 \Term{hilbert-functor}{Hilbert 函子}、\Term{picard-functor}{Picard 函子}等\Term{functor}{函子}的\Term{representability}{可表示性}定理中的\Term{purity}{纯性}假设，值得作系统研究。

$2^\circ$）可以设法用\Term{depth}{深度}给出便于运用的必要充分条件，以保证5.10中考虑的\Term{universal}{普遍}\Term{finite}{有限}性条件成立。当 $S$ 是\Term{field}{域}的谱时，由 EGA III 1.4.15 容易得到：必要且充分条件是各 $R^ig_*(F)$（$i\leq n$）\Term{coherent}{凝聚}；借助 VIII 2.3，这确实可以用\Term{depth}{深度}表达。然而在一般情形，须注意仅要求上述条件对所有\Term{fibre}{纤维} $U_s\subset X_s$（$s\in S$）成立是不够的，即使 $n=0$ 也不够。例如取 $X=S$，其中 $S$ 为\Term{discrete-valuation-ring}{离散赋值环}的谱，$U$ 为仅由\Term{generic-point}{一般点}组成的\Term{open-subset}{开集}，$F=\mathcal O_U$。

$3^\circ$）不过，下面有一个充分条件，保证满足5.10的假设：只须 $f$ \Term{flat}{平坦}，且对每个 $s\in S$ 和每个 $x\in Y_s=X_s-U_s$，都有
\[
 \operatorname{prof}\mathcal O_{X_s,x}\geq n+2.
\]
事实上，考虑到引理2.5（参见2.5之后的关系(16)），可得 $g_*(\mathcal O_U)\simeq\mathcal O_X$，且 $0<i\leq n$ 时 $R^ig_*(\mathcal O_U)=0$，这些关系显然在任意\Term{base-change}{基变换} $S'\to S$ 后仍成立。
